% Co-design queries and the Kleisli category of the upper set monad
% Working note for the ACT4E book, 24 November 2026.
%
% Standalone file: compile with pdflatex (three runs for the references and the table of contents).
% It uses only standard packages and none of the macro files of the book.
\documentclass[11pt,a4paper]{article}

\usepackage[T1]{fontenc}
\usepackage[utf8]{inputenc}
\usepackage{lmodern}
\usepackage[a4paper,hmargin=25mm,vmargin=24mm]{geometry}
\usepackage{amsmath,amssymb,amsthm,mathtools}
\usepackage{stmaryrd}
\usepackage{microtype}
\usepackage{enumitem}
\usepackage{booktabs}
\usepackage{array}
\usepackage{tikz-cd}
\usepackage{graphicx}
\PassOptionsToPackage{hyphens}{url}
\usepackage[hidelinks]{hyperref}
\usepackage[capitalise,noabbrev]{cleveref}
\crefname{equation}{}{}
\Crefname{equation}{Equation}{Equations}
\creflabelformat{equation}{(#2#1#3)}

\theoremstyle{plain}
\newtheorem{theorem}{Theorem}[section]
\newtheorem{proposition}[theorem]{Proposition}
\newtheorem{lemma}[theorem]{Lemma}
\newtheorem{corollary}[theorem]{Corollary}
\theoremstyle{definition}
\newtheorem{definition}[theorem]{Definition}
\newtheorem{example}[theorem]{Example}
\theoremstyle{remark}
\newtheorem{remark}[theorem]{Remark}

% ---------------------------------------------------------------- macros
% A wide check accent (the book has \widecheck from its font setup; here we build one
% by flipping \widehat, so that the note compiles with a plain TeX installation).
\makeatletter
\providecommand{\widecheck}[1]{\mathpalette\note@widecheck{#1}}
\newcommand{\note@widecheck}[2]{%
    \begingroup
    \setbox\z@\hbox{$\m@th#1#2$}%
    \setbox\tw@\hbox{$\m@th#1\widehat{\vrule width\wd\z@ height\z@ depth\z@}$}%
    \setbox4\hbox{$\m@th#1\widehat{#2}$}%
    \rlap{\smash{\raise\dimexpr\ht\z@+\ht\tw@+0.4ex\relax
        \hbox to\wd\z@{\hss\scalebox{1}[-1]{\copy\tw@}\hss}}}%
    \vrule width\z@ height\ht4 depth\z@
    \box\z@
    \endgroup
}
\makeatother
% categories
\newcommand{\Cat}[1]{\mathbf{#1}}
\newcommand{\Set}{\Cat{Set}}
\newcommand{\Rel}{\Cat{Rel}}
\newcommand{\Pre}{\Cat{Pre}}
\newcommand{\Pos}{\Cat{Pos}}
\newcommand{\mRel}{\Cat{mRel}}
\newcommand{\Feas}{\Cat{Feas}}
\newcommand{\DP}{\Cat{DP}}
\newcommand{\BoolProf}{\Cat{Bool}\text{-}\Cat{Prof}}
\newcommand{\Bool}{\Cat{Bool}}
% monads
\newcommand{\U}{\mathsf{U}}
\newcommand{\Lo}{\mathsf{L}}
\newcommand{\pow}{\mathcal{P}}
\newcommand{\An}{\mathcal{A}}
\newcommand{\SetP}{\Set_{\pow}}
\newcommand{\PreU}{\Pre_{\U}}
\newcommand{\PosU}{\Pos_{\U}}
\newcommand{\PreL}{\Pre_{\Lo}}
\newcommand{\PosL}{\Pos_{\Lo}}
\newcommand{\PosA}{\Pos^{\mathrm{wf}}_{\An}}
\newcommand{\PosUwf}{\Pos^{\mathrm{wf}}_{\U}}
% tuple ("strict") versions; the double angle brackets are built from single ones
% in case the font setup does not provide \llangle and \rrangle
\providecommand{\llangle}{\mathopen{\langle\mkern-3.5mu\langle}}
\providecommand{\rrangle}{\mathclose{\rangle\mkern-3.5mu\rangle}}
\newcommand{\tL}[1]{\llangle #1 \rrangle}
% composition, concatenation
\newcommand{\then}{\mathbin{\fatsemi}}
\newcommand{\thenU}{\mathbin{\fatsemi_{\U}}}
\newcommand{\thenL}{\mathbin{\fatsemi_{\Lo}}}
\newcommand{\thenP}{\mathbin{\fatsemi_{\pow}}}
\newcommand{\thenT}{\mathbin{\fatsemi_{T}}}
\newcommand{\thenA}{\mathbin{\fatsemi_{\An}}}
\newcommand{\tc}{\mathbin{\fatsemi_{\!\scriptscriptstyle\langle}}}
\newcommand{\tup}[1]{\langle #1 \rangle}
\newcommand{\etup}{\langle\,\rangle}
% arrows
\newcommand{\profto}{\mathrel{\ooalign{\hfil$\mapstochar\mkern 5mu$\hfil\cr$\to$\cr}}}
\newcommand{\rel}[3]{#1 \mathrel{#2} #3}
% order
\newcommand{\up}{\mathord{\uparrow}}
\newcommand{\dn}{\mathord{\downarrow}}
\newcommand{\op}{^{\mathrm{op}}}
\DeclareMathOperator{\Min}{Min}
\DeclareMathOperator{\Max}{Max}
\DeclareMathOperator{\Tr}{Tr}
\newcommand{\id}{\mathrm{id}}
\newcommand{\N}{\mathbb{N}}
\newcommand{\Z}{\mathbb{Z}}
\newcommand{\R}{\mathbb{R}}
\newcommand{\Rp}{\mathbb{R}_{\geq 0}}
\newcommand{\Fix}{\mathsf{FixFunMinRes}}
\newcommand{\FixR}{\mathsf{FixResMaxFun}}
\newcommand{\comp}[1]{#1_{*}}
\newcommand{\lbl}[1]{\nolinkurl{#1}}
\newcommand{\file}[1]{\nolinkurl{#1}}
\newcommand{\defeq}{\mathrel{:=}}

\title{Co-design queries and the Kleisli category\\ of the upper set monad}
\author{Working note for the ACT4E book:\\ series, parallel, and feedback composition for the query\\ ``fix the functionality, minimize the required resources''}
\date{24 November 2026}

\begin{document}
\maketitle

\tableofcontents

\section{Overview}

A component in co-design is a monotone relation $M \colon F \profto R$ between a preorder $F$ of functionalities and a preorder $R$ of requirements (a \emph{feasibility relation}), or, equivalently, a design problem $d \colon F\op \times R \to \Bool$.
The query ``fix the functionality, minimize the required resources'' ($\Fix$) asks, for a given $f \in F$, for the set
\begin{equation*}
    \widehat{M}(f) = \{\, r \in R \mid \rel{f}{M}{r} \,\}
\end{equation*}
of feasible requirements, or for its minimal elements.
The dual query ``fix the requirement, maximize the functionality'' ($\FixR$) asks, for a given $r \in R$, for the set $\widecheck{M}(r) = \{ f \in F \mid \rel{f}{M}{r} \}$, or for its maximal elements.

This note studies the passage $M \mapsto \widehat{M}$ as a functor, and shows that it is compatible with the three ways of connecting components: in series, in parallel, and with feedback.
Here are the results, in the order of the questions which they answer.

\begin{enumerate}[leftmargin=2em]
    \item \emph{The upper set monad and its Kleisli category} (\cref{sec:upper-set-monad}).
          Write $\U P$ for the set of upper sets of a preorder $P$, ordered by \emph{reverse} inclusion.
          Then $\U$, with unit $p \mapsto \up p$ and multiplication given by union, is a monad on the category $\Pre$ of preorders and monotone functions, and it restricts to a monad on $\Pos$ (\cref{thm:U-monad}).
          Its Kleisli category $\PreU$ is the expected one: a morphism $P \to Q$ is a monotone function $f \colon P \to \U Q$, the composite of $f$ and $g$ sends $p$ to $\bigcup_{q \in f(p)} g(q)$, and the identity sends $p$ to $\up p$ (\cref{prop:kleisli-U}).
          The reverse inclusion order is essential: with the inclusion order, the unit $p \mapsto \up p$ would not be monotone (\cref{rem:order}).
    \item \emph{Upper sets and antichains} (\cref{sec:antichains}).
          For a preorder $P$, upper closure is a bijection from the antichains of $P$ to the upper sets of $P$ if and only if $P$ is a poset in which every non-empty subset has a minimal element, that is, a well-founded poset (\cref{thm:antichains-wf}); the inverse is then $\Min$.
          This covers all finite posets and, for instance, $\N^k$.
          All antichains are moreover finite if and only if $P$ is a well partial order (\cref{prop:wpo}).
          Without well-foundedness, antichains still correspond to the upper sets which are generated by their minimal elements; an example are the topologically closed upper sets of $\Rp^n$ (\cref{exa:closed-upper-sets}).
    \item \emph{Relations are Kleisli morphisms} (\cref{sec:warm-up,sec:mrel-kleisli}).
          The category $\Rel$ is isomorphic to the Kleisli category $\SetP$ of the powerset monad (\cref{thm:rel-kleisli}), and the category $\mRel$ of preorders and monotone relations is isomorphic to $\PreU$, by the functor $\Phi$ which is the identity on objects and sends $M$ to $\widehat{M}$ (\cref{thm:mrel-kleisli}).
          Restricted to posets, this is an isomorphism $\Feas \cong \PosU$, and hence $\DP \cong \PosU$.
          Dually, $M \mapsto \widecheck{M}$ is an isomorphism from $\mRel\op$ to the Kleisli category of the lower set monad (\cref{thm:mrel-kleisli-lower}).
    \item \emph{Parallel composition} (\cref{sec:parallel}).
          The monoidal product of $\PreU$ which corresponds to the product of monotone relations is $(f \otimes g)(p, q) = f(p) \times g(q)$, with the coherence isomorphisms $\tup{\tup{p,q},r} \mapsto \up \tup{p, \tup{q,r}}$, and so on.
          With these structures, $\Phi$ is an isomorphism of symmetric monoidal categories (\cref{thm:PreU-monoidal}).
    \item \emph{Strict versions} (\cref{sec:strict}).
          The same holds for the tuple categories of the book: $\Phi$ restricts to an isomorphism of symmetric strict monoidal categories $\tL{\mRel} \to \tL{\PreU}$, where the monoidal product of $\tL{\PreU}$ concatenates tuples, $(f \otimes g)(a \tc c) = \{ b \tc d \mid b \in f(a),\ d \in g(c) \}$ (\cref{thm:strict-monoidal}).
    \item \emph{Feedback} (\cref{sec:feedback}).
          The formula
          \begin{equation*}
              \Tr^{C}_{A,B}(h)(a) = \{\, b \in B \mid \exists\, c \in C \colon\ b \tc c \in h(a \tc c) \,\}
          \end{equation*}
          defines a trace on $\tL{\PreU}$, and $\Phi \colon \tL{\mRel} \to \tL{\PreU}$ is an isomorphism of traced symmetric strict monoidal categories (\cref{thm:trace}).
          It restricts to $\tL{\Feas} \to \tL{\PosU}$, and composed with the passage from design problems to feasibility relations it is the functor $\Fix \colon \tL{\DP} \to \tL{\PosU}$ (\cref{cor:fixfunminres}).
\end{enumerate}

\Cref{tab:dictionary} summarizes the resulting dictionary.
Its last column describes the same operations in terms of antichains of minimal requirements, for posets in which this is possible (\cref{prop:antichain-category,prop:antichain-product,prop:trace-antichains}).

\begin{table}[ht]
    \centering
    \small
    \renewcommand{\arraystretch}{1.5}
    \begin{tabular}{@{}p{14mm}p{40mm}p{45mm}p{47mm}@{}}
        \toprule
                 & feasibility relations $M$, $N$ & queries $f = \widehat{M}$, $g = \widehat{N}$ (upper sets) & queries $h = \Min \widehat{M}$, $k = \Min \widehat{N}$ (antichains) \\
        \midrule
        identity & $a \leq a'$ & $a \mapsto \up a$ & $a \mapsto \{a\}$ \\
        series   & $\exists\, b \colon \rel{a}{M}{b}$ and $\rel{b}{N}{c}$ & $a \mapsto \bigcup_{b \in f(a)} g(b)$ & $a \mapsto \Min \bigcup_{b \in h(a)} k(b)$ \\
        parallel & $\rel{a}{M}{b}$ and $\rel{c}{N}{d}$ & $a \tc c \mapsto f(a) \tc g(c)$ & $a \tc c \mapsto h(a) \tc k(c)$ \\
        feedback & $\exists\, c \colon \rel{(a \tc c)}{M}{(b \tc c)}$ & $a \mapsto \{ b \mid \exists\, c \colon b \tc c \in f(a \tc c) \}$ & $a \mapsto \Min \{ b \mid \exists\, c \colon b \tc c \in h(a \tc c) \}$ \\
        \bottomrule
    \end{tabular}
    \caption{Design problems and their queries. Here $S \tc T$ denotes the set of all concatenations $s \tc t$ with $s \in S$ and $t \in T$.}
    \label{tab:dictionary}
\end{table}

\paragraph{How the proofs are organized.}
All verifications of axioms are carried out once, in the language of relations, where they are short element chases (\cref{sec:parallel} and \cref{app:mrel-strict,app:mrel-trace}).
The isomorphism $\Phi$ then transports the structures to the Kleisli category (\cref{lem:transport}), and what remains is to compute the transported structures explicitly.
In this way the Kleisli category needs no separate proofs of the pentagon, hexagon, or trace axioms.
\Cref{sec:examples} contains worked computations, and \cref{sec:integration} suggests where the material could go in the book.

\paragraph{Machine check.}
All statements of this note which make sense for finite preorders were also tested by brute force on several thousand randomly generated small preorders and monotone relations: the monad laws, the Kleisli composition formula, functoriality and bijectivity of $\Phi$, all axioms of a symmetric strict monoidal category and of a trace for $\tL{\mRel}$, the compatibility of $\Phi$ with products, braidings, and traces, the dual statements for lower sets, and the antichain formulas.
The numbers in \cref{sec:examples} were computed in the same way.
The script is the file \file{codesign-queries-upper-set-monad-check.py}, which accompanies this note.
This is no substitute for the proofs, but it guards against misprints in the formulas.

\section{Conventions and recollections}
\label{sec:conventions}

References of the form \lbl{def:monotone-relation} are labels in the source of the book.

\subsection{Preorders and upper sets}

A \emph{preorder} $P$ is a set, also denoted $P$, with a reflexive and transitive relation $\leq_P$; it is a \emph{poset} if $\leq_P$ is also antisymmetric.
We write $\Pre$ for the category of preorders and monotone functions, and $\Pos$ for its full subcategory of posets.
We compose in diagrammatic order: for $f \colon P \to Q$ and $g \colon Q \to R$, the composite is $f \then g \colon P \to R$, with $(f \then g)(p) = g(f(p))$.

A subset $S \subseteq P$ is an \emph{upper set} if $p \in S$ and $p \leq p'$ imply $p' \in S$, and a \emph{lower set} if $p \in S$ and $p' \leq p$ imply $p' \in S$.
For $X \subseteq P$ we write
\begin{equation*}
    \up X = \{\, p' \in P \mid \exists\, p \in X \colon p \leq p' \,\},
    \qquad
    \dn X = \{\, p' \in P \mid \exists\, p \in X \colon p' \leq p \,\}
\end{equation*}
for its upper and lower closure, and $\up p = \up \{p\}$, $\dn p = \dn \{p\}$.
For a function $f$ and a subset $X$ of its source, $f[X]$ denotes the image of $X$.

\begin{lemma}
    \label{lem:upper-closure}
    Let $P$ and $Q$ be preorders.
    \begin{enumerate}[label=(\alph*)]
        \item For $X \subseteq P$, the set $\up X$ is an upper set which includes $X$, and it is included in every upper set which includes $X$.
              In particular, $\up S = S$ for every upper set $S$.
        \item Unions of upper sets are upper sets, and $\up \bigcup_{i} X_i = \bigcup_{i} \up X_i$ for every family of subsets $X_i \subseteq P$.
        \item For $p, p' \in P$, we have $p \leq p'$ if and only if $\up p \supseteq \up p'$.
        \item For a monotone function $f \colon P \to Q$ and $X \subseteq P$, we have $\up f[\up X] = \up f[X]$.
        \item For upper sets $S \subseteq P$ and $T \subseteq Q$, the set $S \times T$ is an upper set of the product preorder $P \times Q$, and $\up X \times \up Y = \up (X \times Y)$ for all $X \subseteq P$ and $Y \subseteq Q$.
    \end{enumerate}
\end{lemma}
\begin{proof}
    (a) is \lbl{lem:upper-closure-is-upper-set} and \lbl{lem:upper-closure-smallest} of the book, and (b) is immediate from the definitions.

    (c) If $p \leq p'$ and $p' \leq p''$, then $p \leq p''$ by transitivity, so $\up p' \subseteq \up p$.
    Conversely, $p' \in \up p'$ by reflexivity, so $\up p' \subseteq \up p$ gives $p \leq p'$.

    (d) The inclusion $\supseteq$ holds because $X \subseteq \up X$.
    For $\subseteq$, let $q \in \up f[\up X]$.
    Then there are $p \in X$ and $p' \in P$ with $p \leq p'$ and $f(p') \leq q$.
    Since $f$ is monotone, $f(p) \leq f(p') \leq q$, so $q \in \up f[X]$.

    (e) In $P \times Q$ we have $\tup{p,q} \leq \tup{p',q'}$ if and only if $p \leq p'$ and $q \leq q'$.
    So if $\tup{p,q} \in S \times T$ and $\tup{p,q} \leq \tup{p',q'}$, then $p' \in S$ and $q' \in T$.
    The second claim is immediate from the same description of the order.
\end{proof}

\subsection{Relations and monotone relations}

A relation $R \colon X \to Y$ between sets is a subset $R \subseteq X \times Y$; we write $\rel{x}{R}{y}$ for $\tup{x,y} \in R$.
The composite of $R \colon X \to Y$ and $S \colon Y \to Z$ is
\begin{equation}
    \label{eq:rel-composition}
    R \then S = \{\, \tup{x,z} \mid \exists\, y \in Y \colon \rel{x}{R}{y} \text{ and } \rel{y}{S}{z} \,\},
\end{equation}
and the identity of $X$ is the equality relation.
This is the category $\Rel$.

A \emph{monotone relation} $M \colon P \profto Q$ between preorders (\lbl{def:monotone-relation}) is a relation $M \subseteq P \times Q$ which is
\begin{align}
    \label{eq:closed-source}
    \text{closed downward in the source:} & \quad p' \leq p \text{ and } \rel{p}{M}{q} \ \text{ imply }\ \rel{p'}{M}{q}, \\
    \label{eq:closed-target}
    \text{closed upward in the target:}   & \quad \rel{p}{M}{q} \text{ and } q \leq q' \ \text{ imply }\ \rel{p}{M}{q'}.
\end{align}
Monotone relations compose as relations (\lbl{lem:monrel-composition}), and the preorder relation $\leq_P \colon P \profto P$ is neutral for this composition (\lbl{lem:monrel-leq-neutral}).
So there is a category $\mRel$ whose objects are preorders, whose morphisms are monotone relations, and in which the identity morphism of $P$ is $\leq_P$.
Its full subcategory on the posets is the category $\Feas$ of feasibility relations (\lbl{def:Feas}).

A monotone relation $M \colon P \profto Q$ carries the same information as the Boolean profunctor $\phi_M \colon P\op \times Q \to \Bool$ with $\phi_M(p, q) = \top$ if and only if $\rel{p}{M}{q}$ (\lbl{lem:monrel-as-bool-map}), and all the operations on monotone relations which we use have a counterpart for Boolean profunctors, obtained by replacing ``there exists'' by a join and ``and'' by a meet.
So $\mRel$ is isomorphic to the category $\BoolProf$ of preorders and Boolean profunctors, and $\Feas$ is isomorphic to the category $\DP$ of posets and design problems (\lbl{def:DP}).
We work with monotone relations, and translate to design problems only in the final statements.

Every monotone function $f \colon P \to Q$ has a \emph{companion} (\lbl{def:monrel-companion-conjoint}), the monotone relation
\begin{equation}
    \label{eq:companion}
    \comp{f} \colon P \profto Q,
    \qquad
    \rel{p}{\comp{f}}{q} \iff f(p) \leq q.
\end{equation}
We have $\comp{(\id_P)} = {\leq_P}$ and $\comp{(f \then g)} = \comp{f} \then \comp{g}$ (\lbl{exa:companion-identity}, \lbl{lem:companion-conjoint-composition}).
In other words, taking companions is a functor $\comp{(-)} \colon \Pre \to \mRel$ which is the identity on objects.

\subsection{Monads and their Kleisli categories}

A \emph{monad} $\tup{T, \mu, \eta}$ on a category $\Cat{C}$ (\lbl{def:monad}) consists of a functor $T \colon \Cat{C} \to \Cat{C}$ and natural transformations $\mu \colon TT \Rightarrow T$ and $\eta \colon \id_{\Cat{C}} \Rightarrow T$ such that, for every object $A$,
\begin{equation}
    \label{eq:monad-laws}
    \mu_{TA} \then \mu_A = T\mu_A \then \mu_A,
    \qquad
    \eta_{TA} \then \mu_A = \id_{TA} = T\eta_A \then \mu_A.
\end{equation}
Its \emph{Kleisli category} $\Cat{C}_T$ (\lbl{def:kleisli-category}) has the same objects as $\Cat{C}$; a morphism $A \to B$ of $\Cat{C}_T$ is a morphism $f \colon A \to TB$ of $\Cat{C}$; the composite of $f \colon A \to TB$ and $g \colon B \to TC$ is
\begin{equation}
    \label{eq:kleisli-composition}
    f \thenT g \defeq f \then Tg \then \mu_C \colon A \to TC;
\end{equation}
and the identity morphism of $A$ is $\eta_A$.
There is a functor $J_T \colon \Cat{C} \to \Cat{C}_T$ which is the identity on objects and sends $u \colon A \to B$ to $u \then \eta_B$.
(That $J_T$ preserves composition follows from the naturality of $\eta$ and the unit laws: $(u \then \eta) \thenT (v \then \eta) = u \then \eta \then T v \then T\eta \then \mu = u \then v \then \eta$.)

\section{Warm-up: relations and the powerset monad}
\label{sec:warm-up}

The \emph{powerset monad} $\tup{\pow, \mu, \eta}$ on $\Set$ sends a set $X$ to its set of subsets $\pow X$ and a function $u \colon X \to Y$ to the image function $\pow u \colon \pow X \to \pow Y$, $S \mapsto u[S]$; its unit and multiplication are
\begin{equation*}
    \eta_X \colon X \to \pow X,\ x \mapsto \{x\},
    \qquad \qquad
    \mu_X \colon \pow \pow X \to \pow X,\ \mathcal{S} \mapsto \textstyle\bigcup \mathcal{S}.
\end{equation*}

\begin{lemma}
    \label{lem:powerset-monad}
    $\tup{\pow, \mu, \eta}$ is a monad on $\Set$.
    In its Kleisli category $\SetP$, the composite of $f \colon X \to \pow Y$ and $g \colon Y \to \pow Z$ is given by
    \begin{equation}
        \label{eq:kleisli-powerset}
        (f \thenP g)(x) = \bigcup_{y \in f(x)} g(y).
    \end{equation}
\end{lemma}
\begin{proof}
    $\pow$ is a functor, since $\id[S] = S$ and $v[u[S]] = (u \then v)[S]$.
    Naturality of $\eta$: $u[\{x\}] = \{u(x)\}$.
    Naturality of $\mu$: for $\mathcal{S} \in \pow\pow X$, both $\pow u (\mu_X(\mathcal{S}))$ and $\mu_Y(\pow\pow u(\mathcal{S}))$ are equal to $\bigcup_{S \in \mathcal{S}} u[S]$, because images commute with unions.
    Unit laws: for $S \in \pow X$, we have $\mu_X(\eta_{\pow X}(S)) = \bigcup \{S\} = S$ and $\mu_X(\pow \eta_X(S)) = \bigcup_{x \in S} \{x\} = S$.
    Associativity: for $\mathfrak{S} \in \pow\pow\pow X$, both $\mu_X(\mu_{\pow X}(\mathfrak{S}))$ and $\mu_X(\pow\mu_X(\mathfrak{S}))$ are equal to the set of all $x$ for which there are $S \in \mathcal{S} \in \mathfrak{S}$ with $x \in S$.
    Finally, by \cref{eq:kleisli-composition}, $(f \thenP g)(x) = \mu_Z(g[f(x)]) = \bigcup \{ g(y) \mid y \in f(x) \}$.
\end{proof}

\begin{theorem}
    \label{thm:rel-kleisli}
    There is an isomorphism of categories $\Gamma \colon \Rel \to \SetP$ which is the identity on objects and which sends a relation $R \colon X \to Y$ to the function
    \begin{equation*}
        \widehat{R} \colon X \to \pow Y,
        \qquad
        \widehat{R}(x) = \{\, y \in Y \mid \rel{x}{R}{y} \,\}.
    \end{equation*}
    Its inverse sends a function $f \colon X \to \pow Y$ to the relation $\{ \tup{x,y} \mid y \in f(x) \}$.
\end{theorem}
\begin{proof}
    \emph{Bijection on hom-sets.}
    For a relation $R$ we have $\tup{x,y} \in R$ if and only if $y \in \widehat{R}(x)$, so the relation associated with $\widehat{R}$ is $R$.
    For a function $f \colon X \to \pow Y$ and the relation $R_f = \{ \tup{x,y} \mid y \in f(x) \}$, we have $\widehat{R_f}(x) = \{ y \mid y \in f(x) \} = f(x)$.

    \emph{Identities.}
    The identity of $X$ in $\Rel$ is equality, and $\{ y \mid x = y \} = \{x\} = \eta_X(x)$.

    \emph{Composition.}
    For $R \colon X \to Y$ and $S \colon Y \to Z$, by \cref{eq:rel-composition} and \cref{eq:kleisli-powerset},
    \begin{equation*}
        \widehat{R \then S}(x)
        = \{\, z \mid \exists\, y \colon y \in \widehat{R}(x) \text{ and } z \in \widehat{S}(y) \,\}
        = \bigcup_{y \in \widehat{R}(x)} \widehat{S}(y)
        = (\widehat{R} \thenP \widehat{S})(x).
    \end{equation*}
    So $\Gamma$ is a functor which is bijective on objects and on each hom-set, and therefore an isomorphism of categories; its inverse is the functor described in the statement.
\end{proof}

This is the graded exercise \lbl{ex:HwkRelKleisli} of the book.
Under $\Gamma$, the graph of a function $u \colon X \to Y$ corresponds to $J_{\pow}(u) = u \then \eta_Y$, the function $x \mapsto \{u(x)\}$.

\section{The upper set monad}
\label{sec:upper-set-monad}

\begin{definition}[Upper set monad]
    \label{def:U}
    For a preorder $P$, let $\U P$ be the set of all upper sets of $P$, ordered by reverse inclusion:
    \begin{equation}
        \label{eq:order-UP}
        S \leq_{\U P} T \quad \text{if and only if} \quad S \supseteq T.
    \end{equation}
    The \emph{upper set monad} $\tup{\U, \mu, \eta}$ on $\Pre$ is given by:
    \begin{enumerate}
        \item on objects, $P \mapsto \U P$;
        \item on morphisms, for a monotone function $u \colon P \to Q$,
              \begin{equation*}
                  \U u \colon \U P \to \U Q, \qquad S \mapsto \up u[S];
              \end{equation*}
        \item the unit $\eta_P \colon P \to \U P$, $p \mapsto \up p$;
        \item the multiplication $\mu_P \colon \U\U P \to \U P$, $\mathcal{S} \mapsto \bigcup \mathcal{S}$.
    \end{enumerate}
\end{definition}

Note that $\U P$ is a poset, even if $P$ is only a preorder.
An element $\mathcal{S}$ of $\U\U P$ is a set of upper sets of $P$ which is an upper set of $\U P$; by \cref{eq:order-UP} this means that it is closed under passing to smaller upper sets: if $S \in \mathcal{S}$ and $T \in \U P$ with $T \subseteq S$, then $T \in \mathcal{S}$.
For a set $\mathcal{X}$ of upper sets of $P$, its upper closure in $\U P$ is
\begin{equation}
    \label{eq:upper-closure-in-UP}
    \up \mathcal{X} = \{\, T \in \U P \mid T \subseteq S \text{ for some } S \in \mathcal{X} \,\},
    \qquad \text{and} \qquad
    \textstyle\bigcup \up \mathcal{X} = \bigcup \mathcal{X}.
\end{equation}
For the second equation, ``$\supseteq$'' holds because $\mathcal{X} \subseteq \up \mathcal{X}$, and ``$\subseteq$'' holds because every element of $\up\mathcal{X}$ is included in an element of $\mathcal{X}$.

\begin{theorem}
    \label{thm:U-monad}
    $\tup{\U, \mu, \eta}$ is a monad on $\Pre$.
    Since $\U P$ is always a poset, it restricts to a monad on $\Pos$.
\end{theorem}
\begin{proof}
    \emph{$\U$ is a functor.}
    For monotone $u \colon P \to Q$, the set $\up u[S]$ is an upper set of $Q$, and $S \supseteq S'$ implies $\up u[S] \supseteq \up u[S']$, so $\U u$ is a monotone function $\U P \to \U Q$.
    We have $\U(\id_P)(S) = \up S = S$.
    For monotone $u \colon P \to Q$ and $v \colon Q \to R$, \cref{lem:upper-closure}(d) gives
    \begin{equation*}
        \U v(\U u(S)) = \up v[\up u[S]] = \up v[u[S]] = \U(u \then v)(S).
    \end{equation*}

    \emph{$\eta$ is a natural transformation.}
    By \cref{lem:upper-closure}(c), $\eta_P$ is monotone: $p \leq p'$ implies $\up p \supseteq \up p'$, that is, $\up p \leq_{\U P} \up p'$.
    Naturality: by \cref{lem:upper-closure}(d), $\U u(\eta_P(p)) = \up u[\up p] = \up u(p) = \eta_Q(u(p))$.

    \emph{$\mu$ is a natural transformation.}
    A union of upper sets is an upper set, and $\mathcal{S} \supseteq \mathcal{S}'$ implies $\bigcup \mathcal{S} \supseteq \bigcup \mathcal{S}'$, so $\mu_P$ is a monotone function $\U\U P \to \U P$.
    Naturality: for $\mathcal{S} \in \U\U P$, by \cref{eq:upper-closure-in-UP} and \cref{lem:upper-closure}(b),
    \begin{equation*}
        \mu_Q(\U\U u(\mathcal{S}))
        = \textstyle\bigcup \up \{ \U u(S) \mid S \in \mathcal{S} \}
        = \bigcup_{S \in \mathcal{S}} \up u[S]
        = \up u\bigl[\bigcup \mathcal{S}\bigr]
        = \U u(\mu_P(\mathcal{S})).
    \end{equation*}

    \emph{Unit laws.}
    Let $S \in \U P$.
    By \cref{eq:upper-closure-in-UP},
    \begin{equation*}
        \mu_P(\eta_{\U P}(S)) = \textstyle\bigcup \up \{S\} = S
        \qquad \text{and} \qquad
        \mu_P(\U\eta_P(S)) = \textstyle\bigcup \up \{ \up p \mid p \in S \} = \bigcup_{p \in S} \up p = \up S = S.
    \end{equation*}

    \emph{Associativity.}
    Let $\mathfrak{S} \in \U\U\U P$.
    Then $\mu_P(\mu_{\U P}(\mathfrak{S})) = \bigcup \bigcup \mathfrak{S}$, and by \cref{eq:upper-closure-in-UP},
    \begin{equation*}
        \mu_P(\U\mu_P(\mathfrak{S}))
        = \textstyle\bigcup \up \{ \bigcup \mathcal{S} \mid \mathcal{S} \in \mathfrak{S} \}
        = \bigcup_{\mathcal{S} \in \mathfrak{S}} \bigcup \mathcal{S}
        = \bigcup \bigcup \mathfrak{S}. \qedhere
    \end{equation*}
\end{proof}

\begin{proposition}[The Kleisli category of the upper set monad]
    \label{prop:kleisli-U}
    The Kleisli category $\PreU$ of the upper set monad is the following category.
    \begin{enumerate}
        \item Objects: preorders.
        \item Morphisms $P \to Q$: monotone functions $f \colon P \to \U Q$; that is, functions which assign to each $p \in P$ an upper set $f(p)$ of $Q$, such that
              \begin{equation}
                  \label{eq:kleisli-antitone}
                  p \leq p' \quad \text{implies} \quad f(p) \supseteq f(p').
              \end{equation}
        \item Composition: for $f \colon P \to \U Q$ and $g \colon Q \to \U R$,
              \begin{equation}
                  \label{eq:kleisli-U}
                  (f \thenU g)(p) = \bigcup_{q \in f(p)} g(q).
              \end{equation}
        \item Identity morphisms: $\eta_P \colon P \to \U P$, $p \mapsto \up p$.
    \end{enumerate}
    The functor $J_{\U} \colon \Pre \to \PreU$ sends a monotone function $u \colon P \to Q$ to the function $p \mapsto \up u(p)$.
\end{proposition}
\begin{proof}
    Everything is the definition of the Kleisli category, except \cref{eq:kleisli-U}.
    By \cref{eq:kleisli-composition} and \cref{eq:upper-closure-in-UP},
    \begin{equation*}
        (f \thenU g)(p) = \mu_R(\U g(f(p))) = \textstyle\bigcup \up g[f(p)] = \bigcup g[f(p)] = \bigcup_{q \in f(p)} g(q). \qedhere
    \end{equation*}
\end{proof}

We write $\PosU$ for the full subcategory of $\PreU$ whose objects are the posets; it is the Kleisli category of the restriction of $\U$ to $\Pos$.
This is the category called $\Cat{Pos}_{\U}$ in the draft chapter \file{60_computation/15_query_equiv} (\lbl{def:upos_cat}), where it is defined directly by the formulas of \cref{prop:kleisli-U}.

\begin{remark}[The order on upper sets]
    \label{rem:order}
    The order \cref{eq:order-UP} is forced.
    If we ordered the upper sets by inclusion, the function $p \mapsto \up p$ would reverse the order (\cref{lem:upper-closure}(c)), so it would not be a morphism of $\Pre$, and there would be no unit.
    With reverse inclusion, a larger element of $\U P$ is a smaller set: a harder query has fewer answers.

    The condition \cref{eq:kleisli-antitone} says that a morphism $f \colon P \to Q$ of $\PreU$ is the same thing as a monotone function $P\op \to \tup{\U Q, \subseteq}$.
    This is the form in which the curried design problem $\widehat{d} \colon F\op \to \tup{\U R, \subseteq}$ appears in the book (\lbl{lem:dp-currying-monotone}); the form $F \to \tup{\U R, \supseteq}$ is \lbl{rem:dp-curry-other-form}.
\end{remark}

\begin{remark}
    The poset $\U P$ is a complete lattice in which the meet of a family of upper sets is their union, and $\eta_P \colon P \to \U P$ is the universal monotone function from $P$ to a poset with all meets; so $\U$ is the ``free complete meet-semilattice'' monad.
    We do not use this description.
\end{remark}

\section{Monotone relations are Kleisli morphisms}
\label{sec:mrel-kleisli}

For a monotone relation $M \colon P \profto Q$ and $p \in P$, let
\begin{equation}
    \label{eq:hat}
    \widehat{M}(p) = \{\, q \in Q \mid \rel{p}{M}{q} \,\}.
\end{equation}

\begin{theorem}
    \label{thm:mrel-kleisli}
    There is an isomorphism of categories
    \begin{equation*}
        \Phi \colon \mRel \to \PreU
    \end{equation*}
    which is the identity on objects and which sends a monotone relation $M \colon P \profto Q$ to the function $\widehat{M} \colon P \to \U Q$ of \cref{eq:hat}.
    Its inverse sends a morphism $f \colon P \to \U Q$ of $\PreU$ to the monotone relation
    \begin{equation}
        \label{eq:phi-inverse}
        M_f = \{\, \tup{p,q} \in P \times Q \mid q \in f(p) \,\} \colon P \profto Q.
    \end{equation}
\end{theorem}
\begin{proof}
    \emph{$\Phi$ is well defined.}
    Let $M \colon P \profto Q$ be a monotone relation.
    Each $\widehat{M}(p)$ is an upper set of $Q$: if $q \in \widehat{M}(p)$ and $q \leq q'$, then $\rel{p}{M}{q'}$ by \cref{eq:closed-target}.
    The function $\widehat{M}$ is monotone $P \to \U Q$: if $p \leq p'$ and $q \in \widehat{M}(p')$, then $\rel{p}{M}{q}$ by \cref{eq:closed-source}, so $\widehat{M}(p) \supseteq \widehat{M}(p')$, which is \cref{eq:kleisli-antitone}.

    \emph{The inverse is well defined.}
    Let $f \colon P \to \U Q$ be monotone.
    If $\rel{p}{M_f}{q}$ and $q \leq q'$, then $q' \in f(p)$ because $f(p)$ is an upper set; so $M_f$ is closed upward in the target.
    If $p' \leq p$ and $\rel{p}{M_f}{q}$, then $q \in f(p) \subseteq f(p')$ by \cref{eq:kleisli-antitone}; so $M_f$ is closed downward in the source.

    \emph{Bijection on hom-sets.}
    As in the proof of \cref{thm:rel-kleisli}: $\tup{p,q} \in M_{\widehat{M}}$ if and only if $q \in \widehat{M}(p)$, if and only if $\rel{p}{M}{q}$; and $\widehat{M_f}(p) = \{ q \mid q \in f(p) \} = f(p)$.

    \emph{Identities.}
    The identity of $P$ in $\mRel$ is $\leq_P$, and $\widehat{\leq_P}(p) = \{ p' \mid p \leq p' \} = \up p = \eta_P(p)$.

    \emph{Composition.}
    For $M \colon P \profto Q$ and $N \colon Q \profto R$, by \cref{eq:rel-composition} and \cref{eq:kleisli-U},
    \begin{equation*}
        \widehat{M \then N}(p)
        = \{\, r \mid \exists\, q \colon q \in \widehat{M}(p) \text{ and } r \in \widehat{N}(q) \,\}
        = \bigcup_{q \in \widehat{M}(p)} \widehat{N}(q)
        = (\widehat{M} \thenU \widehat{N})(p).
    \end{equation*}
    Hence $\Phi$ is a functor which is bijective on objects and on hom-sets, so it is an isomorphism of categories.
\end{proof}

So the two conditions in the definition of a monotone relation correspond exactly to the two conditions on a Kleisli morphism: closure upward in the target says that the values are upper sets, and closure downward in the source says that the function is monotone.

\begin{corollary}
    \label{cor:feas-posU}
    $\Phi$ restricts to an isomorphism of categories $\Feas \to \PosU$.
    Composing it with the isomorphism $\DP \to \Feas$, which sends a design problem to its feasible set, we obtain an isomorphism of categories
    \begin{equation*}
        \Fix \colon \DP \to \PosU,
        \qquad
        d \mapsto \widehat{d},
        \qquad
        \widehat{d}(f) = \{\, r \mid d(f^*, r) = \top \,\}.
    \end{equation*}
\end{corollary}
\begin{proof}
    $\Feas$ and $\PosU$ are the full subcategories of $\mRel$ and $\PreU$ on the same class of objects, the posets, and $\Phi$ is the identity on objects.
\end{proof}

\begin{remark}[Literature]
    Both \cref{thm:rel-kleisli} and \cref{thm:mrel-kleisli} are known in substance.
    The first is a standard example of a Kleisli category.
    Monotone relations between posets are the feasibility relations, or $\Bool$-profunctors, of \cite[Chapter 4]{fong-spivak}; the queries are those of \cite{censi}; and the extension of the description of $\Rel$ as a Kleisli category from sets to preorders and posets is treated in \cite{bkpv}.
    What this note adds is a complete and elementary account in the conventions of the book, including parallel composition, the strict versions, and feedback.
\end{remark}

\begin{proposition}[Companions are the free morphisms]
    \label{prop:companion-free}
    For every monotone function $u \colon P \to Q$ we have $\Phi(\comp{u}) = J_{\U}(u)$.
    That is, the triangle of functors
    \begin{equation*}
        \begin{tikzcd}[column sep=large]
            & \Pre \arrow[dl, "\comp{(-)}"'] \arrow[dr, "J_{\U}"] & \\
            \mRel \arrow[rr, "\Phi", "\cong"'] & & \PreU
        \end{tikzcd}
    \end{equation*}
    commutes.
\end{proposition}
\begin{proof}
    All three functors are the identity on objects, and $\widehat{\comp{u}}(p) = \{ q \mid u(p) \leq q \} = \up u(p) = J_{\U}(u)(p)$.
\end{proof}

By \lbl{lem:which-monrels-are-companions}, a monotone relation $M$ is a companion precisely when each $\widehat{M}(p)$ has a least element.
In terms of queries: the morphisms in the image of $J_{\U}$ are the queries which have a least feasible requirement for every functionality.

\begin{remark}[The discrete case]
    \label{rem:discrete}
    Let $D \colon \Set \to \Pre$ be the functor which equips a set with the discrete order.
    The upper sets of $DX$ are all subsets of $X$, so $\U D X$ is $\pow X$ ordered by reverse inclusion, and every function $X \to \pow Y$ is a monotone function $DX \to \U DY$.
    Comparing \cref{eq:kleisli-powerset} with \cref{eq:kleisli-U}, we get a full and faithful functor $\SetP \to \PreU$, which sends $X$ to $DX$ and is the identity on morphisms.
    Likewise there is a full and faithful functor $\Rel \to \mRel$ (\lbl{exa:monrel-discrete}), and the two isomorphisms $\Gamma$ and $\Phi$ are compatible with these functors.
    So \cref{thm:rel-kleisli} is the special case of \cref{thm:mrel-kleisli} for discrete preorders.
\end{remark}

\begin{remark}[Order on hom-sets]
    For monotone relations $M, N \colon P \profto Q$ we have $M \subseteq N$ if and only if $\widehat{M}(p) \subseteq \widehat{N}(p)$ for all $p$, that is, if and only if $\widehat{N} \leq \widehat{M}$ in the point-wise order on functions $P \to \U Q$.
    So $\Phi$ \emph{reverses} the order of the hom-sets: a design problem with more feasible pairs gives a query whose answers are larger sets, that is, smaller elements of $\U Q$.
    Unions and intersections of monotone relations correspond to point-wise unions and intersections of upper sets.
\end{remark}

\section{The dual query: lower sets}
\label{sec:lower}

For a preorder $P$, let $\Lo P$ be the set of lower sets of $P$, ordered by inclusion.
The \emph{lower set monad} $\tup{\Lo, \mu, \eta}$ on $\Pre$ sends a monotone function $u \colon P \to Q$ to $\Lo u \colon S \mapsto \dn u[S]$, and has unit $\eta_P(p) = \dn p$ and multiplication $\mu_P(\mathcal{S}) = \bigcup \mathcal{S}$.

It is obtained from the upper set monad by reversing all orders.
Let $(-)\op \colon \Pre \to \Pre$ be the functor which sends a preorder to its opposite and a monotone function to itself; it is an isomorphism of categories and its own inverse.
The lower sets of $P$ are the upper sets of $P\op$, and inclusion is the opposite of reverse inclusion, so
\begin{equation}
    \label{eq:L-from-U}
    \Lo P = (\U(P\op))\op,
    \qquad
    \Lo u = (\U(u\op))\op,
\end{equation}
and the unit and multiplication of $\Lo$ at $P$ are those of $\U$ at $P\op$, as functions.
Hence the monad laws for $\Lo$ are the monad laws for $\U$, and the Kleisli category $\PreL$ is described as in \cref{prop:kleisli-U}: a morphism $P \to Q$ is a monotone function $g \colon P \to \Lo Q$, that is, a function with values in lower sets such that $p \leq p'$ implies $g(p) \subseteq g(p')$; composition is $(g \thenL g')(p) = \bigcup_{q \in g(p)} g'(q)$; and the identity of $P$ is $p \mapsto \dn p$.

For a monotone relation $M \colon P \profto Q$ and $q \in Q$, let $\widecheck{M}(q) = \{ p \in P \mid \rel{p}{M}{q} \}$.

\begin{theorem}
    \label{thm:mrel-kleisli-lower}
    There is an isomorphism of categories
    \begin{equation*}
        \Psi \colon \mRel\op \to \PreL
    \end{equation*}
    which is the identity on objects and which sends a monotone relation $M \colon P \profto Q$ to the function $\widecheck{M} \colon Q \to \Lo P$.
    In particular,
    \begin{equation}
        \label{eq:check-composition}
        \widecheck{M \then N}(r) = \bigcup_{q \in \widecheck{N}(r)} \widecheck{M}(q)
        \qquad \text{and} \qquad
        \widecheck{\leq_P}(p) = \dn p.
    \end{equation}
    It restricts to an isomorphism $\Feas\op \to \PosL$.
    Composed with the passage from design problems to feasible sets, it gives an isomorphism
    \begin{equation*}
        \FixR \colon \DP\op \to \PosL, \qquad d \mapsto \widecheck{d}.
    \end{equation*}
\end{theorem}
\begin{proof}
    This is the proof of \cref{thm:mrel-kleisli} with the roles of source and target exchanged.
    Each $\widecheck{M}(q)$ is a lower set of $P$ by \cref{eq:closed-source}, and $q \leq q'$ implies $\widecheck{M}(q) \subseteq \widecheck{M}(q')$ by \cref{eq:closed-target}; so $\widecheck{M}$ is a morphism $Q \to P$ of $\PreL$.
    Conversely, for a morphism $g \colon Q \to \Lo P$ of $\PreL$, the relation $\{ \tup{p,q} \mid p \in g(q) \}$ is closed downward in the source because the values of $g$ are lower sets, and closed upward in the target because $g$ is monotone.
    The two assignments are inverse to each other.
    The second equation of \cref{eq:check-composition} is the definition of $\dn p$, and the first one holds because $p \in \widecheck{M \then N}(r)$ means that there is $q$ with $\rel{p}{M}{q}$ and $\rel{q}{N}{r}$, that is, with $q \in \widecheck{N}(r)$ and $p \in \widecheck{M}(q)$.
    The right-hand side is $(\widecheck{N} \thenL \widecheck{M})(r)$, so $\Psi$ reverses the order of composition.
\end{proof}

The query $\FixR$ reads a design problem backwards, from requirements to functionalities; this is why $\Psi$ is contravariant.

\begin{remark}
    \label{rem:U-L}
    By \cref{eq:L-from-U}, a morphism $P \to Q$ of $\PreL$ is the same function as a morphism $P\op \to Q\op$ of $\PreU$, and this is an isomorphism of categories $\PreL \cong \PreU$ which sends $P$ to $P\op$.
    Under this isomorphism, $\Psi$ corresponds to $\Phi$ composed with the isomorphism $\mRel\op \to \mRel$ which sends $P$ to $P\op$ and $M \colon P \profto Q$ to its transpose $Q\op \profto P\op$.

    The draft chapter states that $\FixR$ is a functor $\DP \to \Cat{Pos}_{\Lo}$ (\lbl{lem:confunctor}); since it sends a design problem $F \profto R$ to a morphism $R \to F$, this should read $\DP\op \to \Cat{Pos}_{\Lo}$.
\end{remark}

\section{Upper sets and antichains}
\label{sec:antichains}

In practice one does not want to list all feasible requirements, but only the minimal ones.
This section makes precise when the two descriptions are equivalent, and describes the Kleisli category in terms of antichains.

\subsection{The correspondence}

Let $P$ be a preorder and $X \subseteq P$.
An element $x \in X$ is \emph{minimal} in $X$ if every $y \in X$ with $y \leq x$ satisfies $x \leq y$; if $P$ is a poset, this means that $y \in X$ and $y \leq x$ imply $y = x$.
We write $\Min X$ for the set of minimal elements of $X$.
A subset $A \subseteq P$ is an \emph{antichain} if, for all $a, a' \in A$, $a \leq a'$ implies $a = a'$.
We write $\An P$ for the set of antichains of $P$.
The empty set and all singletons are antichains.

\begin{lemma}
    \label{lem:min-up}
    Let $P$ be a poset and $X \subseteq P$.
    \begin{enumerate}[label=(\alph*)]
        \item $\Min X$ is an antichain.
        \item $\Min \up X = \Min X$.
        \item If $A$ is an antichain, then $\Min \up A = A$.
              In particular, $\up A = \up B$ implies $A = B$ for antichains $A$ and $B$.
    \end{enumerate}
\end{lemma}
\begin{proof}
    (a) If $x, y \in \Min X$ and $x \leq y$, then $x = y$ because $y$ is minimal in $X$.

    (b) Let $x \in \Min X$.
    Then $x \in \up X$.
    If $y \in \up X$ and $y \leq x$, there is $x' \in X$ with $x' \leq y \leq x$; since $x$ is minimal in $X$, $x' = x$, and then $y = x$ by antisymmetry.
    So $x \in \Min \up X$.
    Conversely, let $y \in \Min \up X$.
    There is $x \in X$ with $x \leq y$, and $x \in \up X$, so $x = y$ because $y$ is minimal in $\up X$.
    Hence $y \in X$, and $y$ is minimal in $X$ because $X \subseteq \up X$.

    (c) An antichain satisfies $\Min A = A$, so this follows from (b).
\end{proof}

Parts (a) and (c) are \lbl{lem:Min-is-antichain}, \lbl{lem:up-cl-inj-antichains}, and \lbl{prop:antichains-upper-sets}(1) of the book.
By (c), upper closure is always an injective function $\An P \to \U P$ when $P$ is a poset, and by (b) and (c) its image consists of the upper sets $S$ with $S = \up \Min S$; we say that such an upper set is \emph{generated by its minimal elements}.
(In the book these are called \emph{downward-closed} upper sets, \lbl{def:downward-closed-upperset}.)
So for every poset, $\up$ and $\Min$ are inverse bijections between antichains and upper sets generated by their minimal elements (\lbl{prop:antichains-upper-sets}).
The question is when \emph{every} upper set is of this kind.

\begin{definition}
    A preorder $P$ is \emph{well-founded} if every non-empty subset of $P$ has a minimal element.
\end{definition}

Assuming the axiom of dependent choice, $P$ is well-founded if and only if there is no infinite strictly decreasing sequence $p_0 > p_1 > p_2 > \cdots$ in $P$, where $p > p'$ means $p' \leq p$ and not $p \leq p'$.

\begin{theorem}
    \label{thm:antichains-wf}
    For a preorder $P$, the following are equivalent.
    \begin{enumerate}[label=(\roman*)]
        \item Upper closure $A \mapsto \up A$ is a bijection from the set $\An P$ of antichains of $P$ to the set of upper sets of $P$.
        \item $P$ is a well-founded poset.
    \end{enumerate}
    If these conditions hold, the inverse bijection is $S \mapsto \Min S$.
\end{theorem}
\begin{proof}
    (ii) $\Rightarrow$ (i).
    By \cref{lem:min-up}(c), $\up$ is injective on antichains and $\Min \up A = A$.
    It remains to show that every upper set $S$ satisfies $S = \up \Min S$; then $\up$ is surjective, with inverse $\Min$.
    The inclusion $\up \Min S \subseteq S$ holds because $S$ is an upper set.
    Let $x \in S$.
    The set $X = \{ y \in S \mid y \leq x \}$ is non-empty, so it has a minimal element $m$.
    This $m$ is minimal in $S$: if $y \in S$ and $y \leq m$, then $y \leq x$, so $y \in X$, and hence $y = m$.
    Since $m \leq x$, we get $x \in \up \Min S$.

    (i) $\Rightarrow$ (ii).
    \emph{$P$ is a poset.}
    Suppose $p \leq q$ and $q \leq p$.
    Then $\{p\}$ and $\{q\}$ are antichains with $\up p = \up q$, so $\{p\} = \{q\}$ by injectivity, that is, $p = q$.

    \emph{$P$ is well-founded.}
    Let $X \subseteq P$ be non-empty.
    By surjectivity, $\up X = \up A$ for some antichain $A$, and $A$ is non-empty because $\up X$ is.
    By \cref{lem:min-up}, $A = \Min \up A = \Min \up X = \Min X$.
    So $\Min X$ is non-empty.
\end{proof}

\begin{example}
    \label{exa:wf}
    Every finite poset is well-founded; this is \lbl{lem:finite-upper-sets-dc}.
    The posets $\N$ and $\N^k$ (with the product order) are well-founded, and so is every product of two well-founded posets (\cref{lem:wf-product} below).
    So in the example of paying with coins (\lbl{exa:coins-dp}), where the requirements form the poset $\N \times \N$, every set of feasible requirements is generated by its minimal elements, although the poset is infinite.

    The posets $\Z$, $\tup{\mathbb{Q}_{\geq 0}, \leq}$, and $\tup{\Rp, \leq}$ are not well-founded.
    In $\Rp$, the upper set $\{ x \mid x > 1 \}$ has no minimal element (\lbl{exa:upper-set-no-min}).
    Here the antichains are the empty set and the singletons, which correspond to the upper sets $\emptyset$ and $[x, \infty)$, while the upper sets $(x, \infty)$ are not generated by their minimal elements.
\end{example}

\begin{lemma}
    \label{lem:wf-product}
    If $P$ and $Q$ are well-founded posets, then so is $P \times Q$.
\end{lemma}
\begin{proof}
    Let $X \subseteq P \times Q$ be non-empty.
    Let $p_0$ be a minimal element of $\{ p \mid \tup{p,q} \in X \text{ for some } q \}$, and let $q_0$ be a minimal element of $\{ q \mid \tup{p_0, q} \in X \}$.
    If $\tup{p,q} \in X$ and $\tup{p,q} \leq \tup{p_0,q_0}$, then $p \leq p_0$ gives $p = p_0$, and then $q \leq q_0$ gives $q = q_0$.
    So $\tup{p_0,q_0}$ is minimal in $X$.
\end{proof}

\begin{remark}[Preorders]
    \label{rem:antichains-preorders}
    By \cref{thm:antichains-wf}, the correspondence fails for every preorder which is not a poset: two different elements $p$, $q$ with $p \leq q \leq p$ generate the same upper set.
    This is harmless.
    The upper sets of a preorder $P$ are the preimages of the upper sets of its quotient poset (\file{05_order/20_preorders/72_preorders_to_posets.tex}), so one can apply the theorem to the quotient poset: if it is well-founded, the upper sets of $P$ correspond to the antichains of $P$ \emph{up to replacing elements by equivalent ones}.
\end{remark}

\subsection{Finite antichains}

For computations one would like the antichains to be finite.

\begin{proposition}[Higman]
    \label{prop:wpo}
    For a poset $P$, the following are equivalent.
    \begin{enumerate}[label=(\roman*)]
        \item Every upper set of $P$ is the upper closure of a finite set.
        \item $P$ is well-founded and has no infinite antichain.
        \item Upper closure is a bijection from the set of \emph{finite} antichains of $P$ to the set of upper sets of $P$.
    \end{enumerate}
\end{proposition}
\begin{proof}
    (ii) $\Rightarrow$ (iii): by \cref{thm:antichains-wf}, since all antichains are finite.
    (iii) $\Rightarrow$ (i) is clear.

    (i) $\Rightarrow$ (ii).
    Let $X \subseteq P$ be non-empty, and let $F$ be a finite set with $\up X = \up F$.
    Then $F$ is non-empty and finite, so $\Min F$ is non-empty, and $\Min F = \Min \up F = \Min \up X = \Min X$ by \cref{lem:min-up}(b).
    So $P$ is well-founded.
    If $A$ is an antichain and $\up A = \up F$ with $F$ finite, then $A = \Min \up A = \Min \up F = \Min F \subseteq F$ by \cref{lem:min-up}, so $A$ is finite.
\end{proof}

Posets with these properties are called \emph{well partial orders}, and (i) is called the \emph{finite basis property}; the equivalence goes back to Higman \cite{higman}.
Finite posets and $\N$ are well partial orders, and a product of two well partial orders is again one; for $\N^k$ this is Dickson's lemma (see for instance \cite{dickson-survey}).
So for requirements in $\N^k$ every query has a \emph{finite} antichain of minimal answers.
An infinite set with the discrete order is well-founded but not a well partial order: all its subsets are antichains.

\begin{example}[Closed upper sets of $\Rp^n$]
    \label{exa:closed-upper-sets}
    The poset $\Rp^n$ is not well-founded, but every upper set $S \subseteq \Rp^n$ which is topologically closed is generated by its minimal elements.
    Indeed, let $x \in S$.
    The set $K = \{ y \in S \mid y \leq x \}$ is closed and bounded, hence compact, and it is non-empty.
    So the continuous function $y \mapsto y_1 + \cdots + y_n$ attains a minimum on $K$, at a point $m$.
    If $y \in S$ and $y \leq m$, then $y \in K$ and $y_1 + \cdots + y_n \leq m_1 + \cdots + m_n$, with equality only if $y = m$; by the choice of $m$, equality holds.
    So $m \in \Min S$ and $m \leq x$.

    Therefore $\up$ and $\Min$ are inverse bijections between the closed upper sets of $\Rp^n$ and those antichains whose upper closure is closed.
    Not every antichain is of this kind: for the antichain $A = \{ \tup{t, 1-t} \mid 0 < t < 1 \}$ of $\Rp^2$, the set $\up A$ does not contain its limit point $\tup{0,1}$.
\end{example}

\subsection{Queries in terms of antichains}

Let $P$ be a poset.
Following \lbl{def:antichain-order}, we order the antichains of $P$ by
\begin{equation}
    \label{eq:antichain-order}
    A \leq_{\An P} B
    \quad \text{if and only if} \quad
    \up A \supseteq \up B,
\end{equation}
that is, if and only if for every $b \in B$ there is $a \in A$ with $a \leq b$.
By \cref{lem:min-up}(c) this is a partial order, and by \cref{thm:antichains-wf}, for a well-founded poset $P$ the functions
\begin{equation}
    \label{eq:min-up-iso}
    \Min \colon \U P \to \An P
    \qquad \text{and} \qquad
    \up \colon \An P \to \U P
\end{equation}
are order isomorphisms, inverse to each other.

\begin{proposition}[The category of antichain-valued queries]
    \label{prop:antichain-category}
    The following data define a category $\PosA$.
    \begin{enumerate}
        \item Objects: well-founded posets.
        \item Morphisms $P \to Q$: monotone functions $h \colon P \to \An Q$, that is, functions with values in antichains such that $p \leq p'$ implies $h(p) \leq_{\An Q} h(p')$.
        \item Composition: for $h \colon P \to \An Q$ and $k \colon Q \to \An R$,
              \begin{equation}
                  \label{eq:antichain-composition}
                  (h \thenA k)(p) = \Min \bigcup_{q \in h(p)} k(q).
              \end{equation}
        \item Identity morphisms: $p \mapsto \{p\}$.
    \end{enumerate}
    Let $\PosUwf$ be the full subcategory of $\PosU$ on the well-founded posets.
    There is an isomorphism of categories $\PosUwf \to \PosA$ which is the identity on objects and sends $f \colon P \to \U Q$ to $f \then \Min$, that is, to $p \mapsto \Min f(p)$; its inverse sends $h$ to $p \mapsto \up h(p)$.
\end{proposition}
\begin{proof}
    By \cref{eq:min-up-iso}, composing with $\Min$ is a bijection from monotone functions $P \to \U Q$ to monotone functions $P \to \An Q$, with inverse given by composing with $\up$.
    We transport the composition and the identities of $\PosUwf$ along these bijections; this gives a category which is isomorphic to $\PosUwf$, and we have to show that its composition and identities are the ones in the statement.
    For the identities: $\Min \up p = \{p\}$.
    For the composition, let $h$ and $k$ be as in the statement, and let $f = h \then \up$ and $g = k \then \up$.
    For $p \in P$ we claim
    \begin{equation}
        \label{eq:antichain-composition-proof}
        (f \thenU g)(p)
        = \bigcup_{q \in \up h(p)} \up k(q)
        = \bigcup_{q \in h(p)} \up k(q)
        = \up \bigcup_{q \in h(p)} k(q).
    \end{equation}
    The first equality is \cref{eq:kleisli-U}, and the third is \cref{lem:upper-closure}(b).
    For the second, ``$\supseteq$'' is clear.
    For ``$\subseteq$'', if $q \in \up h(p)$, there is $q_0 \in h(p)$ with $q_0 \leq q$, and $g(q) \subseteq g(q_0)$ by \cref{eq:kleisli-antitone}, because $g$ is a morphism of $\PosU$.
    Applying $\Min$ to \cref{eq:antichain-composition-proof} and using \cref{lem:min-up}(b) gives \cref{eq:antichain-composition}.
\end{proof}

\begin{remark}
    \label{rem:book-series-queries}
    \Cref{eq:antichain-composition} is the formula of \lbl{prop:dp-series-queries} for the minimal requirements function of a series composition, slightly simplified: there the union is taken over the sets $\up k(q)$, which by \cref{lem:min-up}(b) is not necessary.
    In words: collect the minimal requirements $k(q)$ of the second component for each minimal requirement $q$ of the first one, and discard the dominated elements.
    If $Q$ and $R$ are well partial orders, this is a finite computation.
\end{remark}

\begin{remark}
    The upper set monad does \emph{not} restrict to a monad on well-founded posets: for an infinite set $X$ with the discrete order, $\U X$ is the powerset of $X$ ordered by reverse inclusion, which contains infinite strictly decreasing sequences.
    This is the reason why $\PosA$ is defined as a category in its own right, and not as the Kleisli category of an ``antichain monad''.
    On the category of finite posets there is no such problem: there $\An$, with the order \cref{eq:antichain-order}, is a monad which is isomorphic to $\U$.
\end{remark}

\section{Parallel composition}
\label{sec:parallel}

In this section we work with the cartesian product of preorders, which is associative and unital only up to isomorphism; the strict versions follow in \cref{sec:strict}.

\subsection{Monoidal categories and two lemmas on transport}

Recall that a \emph{monoidal category} $\tup{\Cat{C}, \otimes, I, \alpha, \lambda, \rho}$ consists of a category $\Cat{C}$, a functor $\otimes \colon \Cat{C} \times \Cat{C} \to \Cat{C}$, an object $I$, and natural isomorphisms
\begin{equation*}
    \alpha_{A,B,C} \colon (A \otimes B) \otimes C \to A \otimes (B \otimes C),
    \qquad
    \lambda_A \colon I \otimes A \to A,
    \qquad
    \rho_A \colon A \otimes I \to A,
\end{equation*}
which satisfy the pentagon and triangle equations.
It is \emph{symmetric} if it is equipped with a natural isomorphism $\sigma_{A,B} \colon A \otimes B \to B \otimes A$ which satisfies $\sigma_{A,B} \then \sigma_{B,A} = \id$ and the hexagon equation.
It is \emph{strict} if all $\alpha_{A,B,C}$, $\lambda_A$, $\rho_A$ are identity morphisms; this is the notion of \lbl{def:strict-monoidal-cat}.

A functor $G \colon \Cat{C} \to \Cat{D}$ between monoidal categories is \emph{strict monoidal} if it preserves all of this structure on the nose:
\begin{equation}
    \label{eq:strict-monoidal-functor}
    \begin{aligned}
        G(A \otimes B) & = GA \otimes GB, & \qquad G(I) & = I, & \qquad G(u \otimes v) & = Gu \otimes Gv, \\
        G(\alpha) & = \alpha, & G(\lambda) & = \lambda, & G(\rho) & = \rho,
    \end{aligned}
\end{equation}
with the evident indices.
It is \emph{symmetric} if, moreover, $G(\sigma_{A,B}) = \sigma_{GA,GB}$.
An \emph{isomorphism of (symmetric) monoidal categories} is a strict (symmetric) monoidal functor which is an isomorphism of categories; its inverse is then strict (symmetric) monoidal as well.

The category $\Pre$ is symmetric monoidal with the cartesian product: $P \times Q$ carries the product order, $(u \times v)\tup{p,q} = \tup{u(p), v(q)}$, the unit is the one-element preorder $1 = \{*\}$, and
\begin{equation}
    \label{eq:pre-coherence}
    \alpha\tup{\tup{p,q},r} = \tup{p,\tup{q,r}},
    \qquad
    \lambda\tup{*,p} = p,
    \qquad
    \rho\tup{p,*} = p,
    \qquad
    \sigma\tup{p,q} = \tup{q,p}.
\end{equation}

The first lemma says that structure can be moved along an isomorphism of categories.
We state it for functors which are the identity on objects, which is all we need.

\begin{lemma}[Transport of structure]
    \label{lem:transport}
    Let $\Cat{C}$ and $\Cat{D}$ be categories with the same objects, and let $\Phi \colon \Cat{C} \to \Cat{D}$ be an isomorphism of categories which is the identity on objects.
    \begin{enumerate}[label=(\alph*)]
        \item If $\tup{\otimes, I, \alpha, \lambda, \rho}$ is a monoidal structure on $\Cat{C}$, then
              \begin{equation}
                  \label{eq:transport}
                  u \otimes' v \defeq \Phi\bigl(\Phi^{-1}u \otimes \Phi^{-1}v\bigr),
                  \qquad
                  \alpha' \defeq \Phi(\alpha),
                  \quad
                  \lambda' \defeq \Phi(\lambda),
                  \quad
                  \rho' \defeq \Phi(\rho),
              \end{equation}
              together with the product of objects and the unit of $\Cat{C}$, is a monoidal structure on $\Cat{D}$.
              It is strict if the one on $\Cat{C}$ is.
        \item If, moreover, $\sigma$ is a symmetry for $\Cat{C}$, then $\sigma' \defeq \Phi(\sigma)$ is a symmetry for $\Cat{D}$.
        \item If, moreover, $\Cat{C}$ is strict and $\Tr$ is a feedback operator, or a trace, on $\Cat{C}$ (\lbl{def:feedback-cat}, \lbl{def:trace-operator}), then $\Tr'(h) \defeq \Phi(\Tr(\Phi^{-1}h))$ is a feedback operator, respectively a trace, on $\Cat{D}$.
    \end{enumerate}
    In each case, $\Phi$ preserves the structure on the nose, and the structure on $\Cat{D}$ is the only one with this property.
\end{lemma}
\begin{proof}
    The operation $\otimes'$ is the composite functor $\Cat{D} \times \Cat{D} \to \Cat{C} \times \Cat{C} \to \Cat{C} \to \Cat{D}$ of $\Phi^{-1} \times \Phi^{-1}$, $\otimes$, and $\Phi$, so it is a functor, and by construction $\Phi(u_0 \otimes v_0) = \Phi u_0 \otimes' \Phi v_0$ for all morphisms $u_0$, $v_0$ of $\Cat{C}$.
    Every condition in the definitions of a monoidal structure, of a symmetry, and of a feedback operator or trace is of one of the following kinds: an equation between objects built from $\otimes$ and $I$, which is the same statement for $\Cat{C}$ and $\Cat{D}$; the requirement that a structure morphism be invertible, which is preserved by the functor $\Phi$; or an equation between two morphisms which are built from arbitrary morphisms $u, v, \ldots$ and from the structure morphisms by means of composition, $\otimes$, and $\Tr$ (naturality conditions are of this kind).
    Given morphisms $u, v, \ldots$ of $\Cat{D}$, such an equation in $\Cat{D}$ is the image under $\Phi$ of the corresponding equation in $\Cat{C}$ for the morphisms $\Phi^{-1}u, \Phi^{-1}v, \ldots$, because $\Phi$ preserves composition and identities and turns $\otimes$, $\Tr$, and the structure morphisms of $\Cat{C}$ into $\otimes'$, $\Tr'$, and the structure morphisms of $\Cat{D}$.
    So all conditions hold in $\Cat{D}$.
    The formulas for the structure on $\Cat{D}$ are forced by the requirement that $\Phi$ preserve the structure, which gives uniqueness.
\end{proof}

The second lemma produces a monoidal structure from one on a smaller category.

\begin{lemma}[Extension along an identity-on-objects functor]
    \label{lem:extension}
    Let $\tup{\Cat{C}, \otimes, I, \alpha, \lambda, \rho}$ be a monoidal category, let $\Cat{D}$ be a category with the same objects, and let $J \colon \Cat{C} \to \Cat{D}$ be a functor which is the identity on objects.
    Let $\boxtimes \colon \Cat{D} \times \Cat{D} \to \Cat{D}$ be a functor which agrees with $\otimes$ on objects and satisfies $J(u \otimes v) = Ju \boxtimes Jv$ for all morphisms $u$, $v$ of $\Cat{C}$.
    Suppose that the families $J\alpha$, $J\lambda$, and $J\rho$ are natural with respect to all morphisms of $\Cat{D}$.
    Then $\tup{\Cat{D}, \boxtimes, I, J\alpha, J\lambda, J\rho}$ is a monoidal category, and $J$ is a strict monoidal functor.
    If, moreover, $\sigma$ is a symmetry for $\Cat{C}$ and the family $J\sigma$ is natural with respect to all morphisms of $\Cat{D}$, then $J\sigma$ is a symmetry for $\Cat{D}$, and $J$ is symmetric.
\end{lemma}
\begin{proof}
    The morphisms $J\alpha_{A,B,C}$, $J\lambda_A$, $J\rho_A$, $J\sigma_{A,B}$ are isomorphisms because functors preserve isomorphisms, and they are natural by assumption.
    Each of the remaining conditions (pentagon, triangle, hexagon, $\sigma_{A,B} \then \sigma_{B,A} = \id$) is an equation between composites of morphisms of the form $x \otimes y$, where $x$ and $y$ are identities or components of $\alpha$, $\lambda$, $\rho$, $\sigma$.
    Applying $J$ to such an equation in $\Cat{C}$, and using $J(x \otimes y) = Jx \boxtimes Jy$ and $J(\id) = \id$, gives the corresponding equation in $\Cat{D}$.
    That $J$ is strict (symmetric) monoidal is then true by construction.
\end{proof}

\subsection{The monoidal category of monotone relations}

For monotone relations $M \colon P \profto P'$ and $N \colon Q \profto Q'$, their product (\lbl{def:monrel-product}) is the monotone relation $M \times N \colon P \times Q \profto P' \times Q'$ with
\begin{equation}
    \label{eq:monrel-product}
    \rel{\tup{p,q}}{(M \times N)}{\tup{p',q'}}
    \iff
    \rel{p}{M}{p'} \text{ and } \rel{q}{N}{q'}.
\end{equation}

\begin{lemma}[Sliding]
    \label{lem:sliding}
    Let $u \colon P \to Q$ be a monotone function and $N \colon Q \profto R$ a monotone relation.
    Then
    \begin{equation}
        \label{eq:sliding-left}
        \comp{u} \then N = \{\, \tup{p,r} \mid \rel{u(p)}{N}{r} \,\}.
    \end{equation}
    Let $M \colon P \profto Q$ be a monotone relation and $v \colon Q \to R$ an order isomorphism.
    Then
    \begin{equation}
        \label{eq:sliding-right}
        M \then \comp{v} = \{\, \tup{p,r} \mid \rel{p}{M}{v^{-1}(r)} \,\}.
    \end{equation}
\end{lemma}
\begin{proof}
    By \cref{eq:companion}, $\rel{p}{(\comp{u} \then N)}{r}$ means that there is $q$ with $u(p) \leq q$ and $\rel{q}{N}{r}$.
    If so, then $\rel{u(p)}{N}{r}$ by \cref{eq:closed-source}; conversely, we can take $q = u(p)$.
    Similarly, $\rel{p}{(M \then \comp{v})}{r}$ means that there is $q$ with $\rel{p}{M}{q}$ and $v(q) \leq r$, that is, $q \leq v^{-1}(r)$, since $v$ is an order isomorphism.
    If so, then $\rel{p}{M}{v^{-1}(r)}$ by \cref{eq:closed-target}; conversely, we can take $q = v^{-1}(r)$.
\end{proof}

\begin{proposition}
    \label{prop:mrel-monoidal}
    $\mRel$ is a symmetric monoidal category with the product \cref{eq:monrel-product}, the unit $1$, and the companions $\comp{\alpha}$, $\comp{\lambda}$, $\comp{\rho}$, $\comp{\sigma}$ of the coherence isomorphisms \cref{eq:pre-coherence} of $\Pre$.
    The functor $\comp{(-)} \colon \Pre \to \mRel$ is strict symmetric monoidal.
\end{proposition}
\begin{proof}
    We verify the hypotheses of \cref{lem:extension} for $J = \comp{(-)}$.

    \emph{The product is a functor.}
    $M \times N$ is a monotone relation by \lbl{lem:monrel-product}.
    It preserves identities, ${\leq_P} \times {\leq_Q} = {\leq_{P \times Q}}$, by the definition of the product order.
    It preserves composition: for $M' \colon P' \profto P''$ and $N' \colon Q' \profto Q''$, both $(M \then M') \times (N \then N')$ and $(M \times N) \then (M' \times N')$ relate $\tup{p,q}$ to $\tup{p'',q''}$ precisely when there are $p'$ and $q'$ with $\rel{p}{M}{p'}$, $\rel{p'}{M'}{p''}$, $\rel{q}{N}{q'}$, and $\rel{q'}{N'}{q''}$.

    \emph{Compatibility with $J$.}
    For monotone functions $u$, $v$, both $\comp{(u \times v)}$ and $\comp{u} \times \comp{v}$ relate $\tup{p,q}$ to $\tup{p',q'}$ precisely when $u(p) \leq p'$ and $v(q) \leq q'$.

    \emph{Naturality.}
    Let $M \colon P \profto P'$, $N \colon Q \profto Q'$, and $L \colon R \profto R'$ be monotone relations.
    By \cref{eq:sliding-right} and \cref{eq:sliding-left}, both
    \begin{equation*}
        ((M \times N) \times L) \then \comp{\alpha}
        \qquad \text{and} \qquad
        \comp{\alpha} \then (M \times (N \times L))
    \end{equation*}
    relate $\tup{\tup{p,q},r}$ to $\tup{p',\tup{q',r'}}$ precisely when $\rel{p}{M}{p'}$, $\rel{q}{N}{q'}$, and $\rel{r}{L}{r'}$.
    In the same way, both $({\leq_1} \times M) \then \comp{\lambda}$ and $\comp{\lambda} \then M$ relate $\tup{*,p}$ to $p'$ precisely when $\rel{p}{M}{p'}$, and similarly for $\rho$; and both $(M \times N) \then \comp{\sigma}$ and $\comp{\sigma} \then (N \times M)$ relate $\tup{p,q}$ to $\tup{q',p'}$ precisely when $\rel{p}{M}{p'}$ and $\rel{q}{N}{q'}$.
\end{proof}

Explicitly, the coherence isomorphisms of $\mRel$ are obtained from those of $\Pre$ by replacing equality with the order: for instance, $\comp{\alpha}$ relates $\tup{\tup{p,q},r}$ to $\tup{p',\tup{q',r'}}$ precisely when $p \leq p'$, $q \leq q'$, and $r \leq r'$.

\subsection{The monoidal category of queries}

\begin{theorem}
    \label{thm:PreU-monoidal}
    The following data define a symmetric monoidal structure on $\PreU$.
    \begin{enumerate}
        \item On objects, the monoidal product is the product of preorders, and the unit is $1$.
        \item On morphisms, for $f \colon P \to \U P'$ and $g \colon Q \to \U Q'$,
              \begin{equation}
                  \label{eq:kleisli-product}
                  f \otimes g \colon P \times Q \to \U(P' \times Q'),
                  \qquad
                  (f \otimes g)\tup{p,q} = f(p) \times g(q).
              \end{equation}
        \item The coherence isomorphisms are the images under $J_{\U}$ of those of $\Pre$:
              \begin{equation}
                  \label{eq:kleisli-coherence}
                  \begin{aligned}
                      \alpha^{\U}\tup{\tup{p,q},r} & = \up \tup{p,\tup{q,r}}, &
                      \lambda^{\U}\tup{*,p} & = \up p, \\
                      \sigma^{\U}\tup{p,q} & = \up \tup{q,p}, &
                      \rho^{\U}\tup{p,*} & = \up p.
                  \end{aligned}
              \end{equation}
    \end{enumerate}
    With this structure, $\Phi \colon \mRel \to \PreU$ is an isomorphism of symmetric monoidal categories, and $J_{\U} \colon \Pre \to \PreU$ is a strict symmetric monoidal functor.
    In particular, for all monotone relations $M$ and $N$,
    \begin{equation}
        \label{eq:hat-product}
        \widehat{M \times N}\tup{p,q} = \widehat{M}(p) \times \widehat{N}(q).
    \end{equation}
\end{theorem}
\begin{proof}
    By \cref{lem:transport}, applied to $\Phi$ and to the symmetric monoidal structure of $\mRel$ (\cref{prop:mrel-monoidal}), there is exactly one symmetric monoidal structure on $\PreU$ for which $\Phi$ is an isomorphism of symmetric monoidal categories.
    We compute it.
    The product of objects and the unit are those of $\mRel$.
    For morphisms $f$ and $g$ as in the statement, \cref{eq:transport}, \cref{eq:phi-inverse}, and \cref{eq:monrel-product} give
    \begin{equation*}
        \Phi(M_f \times M_g)\tup{p,q}
        = \{\, \tup{p',q'} \mid p' \in f(p) \text{ and } q' \in g(q) \,\}
        = f(p) \times g(q),
    \end{equation*}
    which is \cref{eq:kleisli-product}; taking $f = \widehat{M}$ and $g = \widehat{N}$ gives \cref{eq:hat-product}.
    The coherence isomorphisms are $\Phi(\comp{\alpha}) = J_{\U}(\alpha)$, and so on, by \cref{prop:companion-free}, and $J_{\U}(u)(x) = \up u(x)$.
    Finally, $J_{\U}$ is the composite of the strict symmetric monoidal functors $\comp{(-)}$ and $\Phi$.
\end{proof}

It is worth making explicit what the theorem contains.
The set $f(p) \times g(q)$ is an upper set of $P' \times Q'$ (\cref{lem:upper-closure}(e)), and $f \otimes g$ is monotone.
The product is a functor: it preserves identities, $\up p \times \up q = \up \tup{p,q}$, and composition,
\begin{equation}
    \label{eq:kleisli-interchange}
    (f \thenU f') \otimes (g \thenU g') = (f \otimes g) \thenU (f' \otimes g'),
\end{equation}
which at $\tup{p,q}$ is the equation $\bigl(\bigcup_{p' \in f(p)} f'(p')\bigr) \times \bigl(\bigcup_{q' \in g(q)} g'(q')\bigr) = \bigcup_{\tup{p',q'} \in f(p) \times g(q)} f'(p') \times g'(q')$.
Naturality of $\alpha^{\U}$ says that, for all morphisms $f$, $g$, $e$ of $\PreU$, the upper sets $(f(p) \times g(q)) \times e(r)$ and $f(p) \times (g(q) \times e(r))$ correspond to each other under the order isomorphism $\alpha$.
All of these facts can also be checked directly; here they follow from the corresponding facts for relations.

\begin{remark}[Where the product comes from]
    \label{rem:monoidal-monad}
    The product \cref{eq:kleisli-product} is an instance of a general construction.
    The functions
    \begin{equation*}
        \nabla_{P,Q} \colon \U P \times \U Q \to \U(P \times Q),
        \qquad
        \tup{S,T} \mapsto S \times T,
    \end{equation*}
    are monotone and natural in $P$ and $Q$, they are compatible with the coherence isomorphisms \cref{eq:pre-coherence}, and they are compatible with the unit and the multiplication of the monad:
    \begin{equation*}
        (\eta_P \times \eta_Q) \then \nabla_{P,Q} = \eta_{P \times Q},
        \qquad
        (\mu_P \times \mu_Q) \then \nabla_{P,Q} = \nabla_{\U P, \U Q} \then \U\nabla_{P,Q} \then \mu_{P \times Q}.
    \end{equation*}
    (The first equation is $\up p \times \up q = \up\tup{p,q}$, and both sides of the second one send $\tup{\mathcal{S},\mathcal{T}}$ to $\bigcup\mathcal{S} \times \bigcup\mathcal{T}$.)
    So $\U$ is a symmetric monoidal monad, also called a commutative monad, on $\tup{\Pre, \times, 1}$.
    For every such monad the Kleisli category inherits a symmetric monoidal structure, with $f \otimes g = (f \times g) \then \nabla$ and with coherence isomorphisms the images of those of the base category; see \cite{kock,power-robinson}.
    This is the structure of \cref{thm:PreU-monoidal}.
    The powerset monad is commutative in the same way, and $\Gamma \colon \Rel \to \SetP$ is an isomorphism of symmetric monoidal categories; by \cref{rem:discrete} this is the discrete case.
\end{remark}

\begin{remark}
    Although the monoidal product of $\PreU$ is the cartesian product on objects, it is \emph{not} the categorical product of $\PreU$, just as the cartesian product of sets is not the categorical product in $\Rel$.
\end{remark}

\subsection{Parallel composition in terms of antichains}

\begin{lemma}
    \label{lem:min-product}
    Let $P$ and $Q$ be posets, $X \subseteq P$, and $Y \subseteq Q$.
    Then $\Min(X \times Y) = \Min X \times \Min Y$.
    In particular, a product of two antichains is an antichain.
\end{lemma}
\begin{proof}
    If $x \in \Min X$ and $y \in \Min Y$, and $\tup{x',y'} \in X \times Y$ with $\tup{x',y'} \leq \tup{x,y}$, then $x' = x$ and $y' = y$.
    Conversely, let $\tup{x,y}$ be minimal in $X \times Y$.
    If $x' \in X$ and $x' \leq x$, then $\tup{x',y} \leq \tup{x,y}$, so $x' = x$; hence $x \in \Min X$, and in the same way $y \in \Min Y$.
\end{proof}

\begin{proposition}
    \label{prop:antichain-product}
    The symmetric monoidal structure of $\PreU$ restricts to the full subcategory $\PosUwf$ on the well-founded posets.
    Under the isomorphism $\PosUwf \cong \PosA$ of \cref{prop:antichain-category}, the monoidal product of $h \colon P \to \An P'$ and $k \colon Q \to \An Q'$ is
    \begin{equation}
        \label{eq:antichain-product}
        (h \otimes k)\tup{p,q} = h(p) \times k(q),
    \end{equation}
    and the coherence isomorphisms are $\tup{\tup{p,q},r} \mapsto \{\tup{p,\tup{q,r}}\}$, and so on.
\end{proposition}
\begin{proof}
    A product of two well-founded posets is well-founded (\cref{lem:wf-product}), and $1$ is well-founded.
    For $f = h \then \up$ and $g = k \then \up$, \cref{lem:upper-closure}(e), \cref{lem:min-up}(b), and \cref{lem:min-product} give
    \begin{equation*}
        \Min\bigl( f(p) \times g(q) \bigr) = \Min \up \bigl( h(p) \times k(q) \bigr) = \Min\bigl( h(p) \times k(q) \bigr) = h(p) \times k(q).
    \end{equation*}
    For the coherence isomorphisms, use $\Min \up x = \{x\}$.
\end{proof}

In words: the minimal requirements of two components in parallel are the pairs of minimal requirements of the two components; no domination check is needed.

\section{Strict versions}
\label{sec:strict}

\subsection{Tuple preorders}

We recall the constructions of the book (\file{27_interconnections/20_parallelism/24_monoidal.tex}).
For preorders $P_1, \ldots, P_n$, with $n \in \N$, the \emph{tuple preorder} $\tL{P_1, \ldots, P_n}$ is the set of all tuples $\tup{p_1, \ldots, p_n}$ with $p_i \in P_i$, ordered entry by entry.
For $n = 0$ it is the preorder $\tL{\,} = \{\etup\}$ whose only element is the empty tuple.
We write $a \tc b$ for the concatenation of two tuples; it is associative, and the empty tuple is neutral for it.
For sets of tuples $S$ and $T$ we write
\begin{equation}
    \label{eq:concat-sets}
    S \tc T = \{\, a \tc b \mid a \in S,\ b \in T \,\}.
\end{equation}

The category $\tL{\Pre}$ is the full subcategory of $\Pre$ on the tuple preorders (\lbl{def:PreL}).
It is a strict monoidal category (\lbl{lem:PreL-strict-monoidal}): the monoidal product of objects concatenates the lists of preorders,
\begin{equation*}
    \tL{P_1, \ldots, P_m} \otimes \tL{Q_1, \ldots, Q_n} = \tL{P_1, \ldots, P_m, Q_1, \ldots, Q_n},
\end{equation*}
so that, as a set, $P \otimes Q = P \tc Q$; the monoidal product of monotone functions is $(u \otimes v)(a \tc c) = u(a) \tc v(c)$; and the unit is $\tL{\,}$.
It is symmetric, with braiding $a \tc b \mapsto b \tc a$ (\lbl{exa:SetL-symmetric}).\footnote{If one of the $P_i$ is empty, then $\tL{P_1, \ldots, P_n}$ is empty and does not determine the list $P_1, \ldots, P_n$. This causes no problems: the monoidal product of an empty object with any object is empty, and every statement below is trivially true when one of the objects involved is empty. A non-empty tuple preorder determines its list.}

\begin{lemma}[Splitting]
    \label{lem:splitting}
    Let $P$ and $Q$ be objects of $\tL{\Pre}$.
    Every element of $P \otimes Q$ can be written as $a \tc b$ with $a \in P$ and $b \in Q$ in exactly one way, and
    \begin{equation*}
        a \tc b \leq a' \tc b' \quad \text{if and only if} \quad a \leq a' \text{ and } b \leq b'.
    \end{equation*}
    Consequently, $S \tc T$ is an upper set of $P \otimes Q$ for all upper sets $S \subseteq P$ and $T \subseteq Q$, and $\up X \tc \up Y = \up (X \tc Y)$ for all subsets $X \subseteq P$ and $Y \subseteq Q$.
\end{lemma}
\begin{proof}
    If $P = \tL{P_1, \ldots, P_m}$ and $Q = \tL{Q_1, \ldots, Q_n}$, a tuple of length $m + n$ is the concatenation of its first $m$ entries and its last $n$ entries, and of no other pair of tuples of lengths $m$ and $n$; and tuples are compared entry by entry.
    The consequences follow as in \cref{lem:upper-closure}(e).
\end{proof}

In other words, $\tup{a,b} \mapsto a \tc b$ is an order isomorphism $P \times Q \to P \otimes Q$.
The strict monoidal product does the same job as the cartesian product, with concatenation instead of pairing.

The categories $\tL{\mRel}$ and $\tL{\Feas}$ of the book are the full subcategories of $\mRel$ on the tuple preorders, and on the tuple preorders which are posets, respectively (\lbl{def:mRelL}, \lbl{def:DPL}).

\begin{proposition}
    \label{prop:mRelL}
    $\tL{\mRel}$ is a symmetric strict monoidal category (\lbl{def:sym-stacking-cat}) with the following structure.
    On objects, the monoidal product and the unit are those of $\tL{\Pre}$.
    For monotone relations $M \colon A \profto B$ and $N \colon C \profto D$,
    \begin{equation}
        \label{eq:mRelL-product}
        M \otimes N = \{\, \tup{a \tc c,\ b \tc d} \mid \rel{a}{M}{b} \text{ and } \rel{c}{N}{d} \,\} \colon A \otimes C \profto B \otimes D,
    \end{equation}
    and the braiding is
    \begin{equation}
        \label{eq:mRelL-braiding}
        \sigma_{A,B} = \{\, \tup{a \tc b,\ b' \tc a'} \mid a \leq a' \text{ and } b \leq b' \,\} \colon A \otimes B \profto B \otimes A.
    \end{equation}
    The structure restricts to $\tL{\Feas}$.
\end{proposition}

This is stated in the book (\lbl{eq:RelL-monoidal-product}, \lbl{exa:RelL-symmetric}).
For completeness, \cref{app:mrel-strict} contains a proof.

\subsection{The strict monoidal category of queries}

\begin{definition}
    \label{def:PreUL}
    $\tL{\PreU}$ is the full subcategory of $\PreU$ whose objects are the tuple preorders, and $\tL{\PosU}$ is the full subcategory of $\PreU$ whose objects are the tuple preorders which are posets.
\end{definition}

So a morphism $A \to B$ of $\tL{\PreU}$ is a monotone function $f \colon A \to \U B$ between tuple preorders, and composition and identities are as in \cref{prop:kleisli-U}.
The isomorphism $\Phi$ of \cref{thm:mrel-kleisli} is the identity on objects, so it restricts to isomorphisms of categories
\begin{equation*}
    \Phi \colon \tL{\mRel} \to \tL{\PreU}
    \qquad \text{and} \qquad
    \Phi \colon \tL{\Feas} \to \tL{\PosU}.
\end{equation*}

\begin{theorem}
    \label{thm:strict-monoidal}
    The following data define a symmetric strict monoidal structure on $\tL{\PreU}$.
    \begin{enumerate}
        \item On objects, the monoidal product and the unit are those of $\tL{\Pre}$.
        \item On morphisms, for $f \colon A \to \U B$ and $g \colon C \to \U D$, the morphism $f \otimes g \colon A \otimes C \to \U(B \otimes D)$ is
              \begin{equation}
                  \label{eq:strict-kleisli-product}
                  (f \otimes g)(a \tc c) = f(a) \tc g(c) = \{\, b \tc d \mid b \in f(a),\ d \in g(c) \,\}.
              \end{equation}
        \item The braiding is
              \begin{equation}
                  \label{eq:strict-kleisli-braiding}
                  \sigma_{A,B} \colon A \otimes B \to \U(B \otimes A),
                  \qquad
                  \sigma_{A,B}(a \tc b) = \up (b \tc a) = \up b \tc \up a.
              \end{equation}
    \end{enumerate}
    With this structure, $\Phi \colon \tL{\mRel} \to \tL{\PreU}$ is an isomorphism of symmetric strict monoidal categories; in particular,
    \begin{equation}
        \label{eq:hat-strict-product}
        \widehat{M \otimes N}(a \tc c) = \widehat{M}(a) \tc \widehat{N}(c).
    \end{equation}
    The structure restricts to $\tL{\PosU}$, and $\Phi \colon \tL{\Feas} \to \tL{\PosU}$ is an isomorphism of symmetric strict monoidal categories.
    Moreover, $J_{\U}$ restricts to a strict symmetric monoidal functor $\tL{\Pre} \to \tL{\PreU}$.
\end{theorem}
\begin{proof}
    By \cref{lem:transport} and \cref{prop:mRelL}, there is exactly one symmetric strict monoidal structure on $\tL{\PreU}$ for which $\Phi$ is an isomorphism of symmetric strict monoidal categories, and we compute it.
    The product of objects and the unit are those of $\tL{\mRel}$, hence of $\tL{\Pre}$.
    For $f$ and $g$ as in the statement, \cref{eq:transport}, \cref{eq:phi-inverse}, and \cref{eq:mRelL-product} give
    \begin{equation*}
        \Phi(M_f \otimes M_g)(a \tc c) = \{\, b \tc d \mid b \in f(a) \text{ and } d \in g(c) \,\} = f(a) \tc g(c).
    \end{equation*}
    By \cref{eq:mRelL-braiding}, $\Phi(\sigma_{A,B})(a \tc b) = \{ b' \tc a' \mid a \leq a',\ b \leq b' \} = \up b \tc \up a$, which is $\up(b \tc a)$ by \cref{lem:splitting}.
    A monoidal product of two posets is a poset, and the unit is a poset, so the structure restricts to $\tL{\PosU}$ and $\tL{\Feas}$.
    For a monotone function $u$ between tuple preorders, $\Phi(\comp{u}) = J_{\U}(u)$ by \cref{prop:companion-free}.
    The functor $\comp{(-)} \colon \tL{\Pre} \to \tL{\mRel}$ is strict symmetric monoidal: both $\comp{(u \otimes v)}$ and $\comp{u} \otimes \comp{v}$ relate $a \tc c$ to $b \tc d$ precisely when $u(a) \leq b$ and $v(c) \leq d$, and the braiding \cref{eq:mRelL-braiding} is the companion of $a \tc b \mapsto b \tc a$.
    Hence so is the composite $J_{\U}$ of $\comp{(-)}$ and $\Phi$.
\end{proof}

Let us spell out what strictness means for queries.
For morphisms $f$, $g$, $e$ of $\tL{\PreU}$, both $(f \otimes g) \otimes e$ and $f \otimes (g \otimes e)$ send $a \tc b \tc c$ to $f(a) \tc g(b) \tc e(c)$.
The identity of the unit $\tL{\,}$ sends $\etup$ to $\up \etup = \{\etup\}$, and so $(f \otimes \id_{\tL{\,}})(a \tc \etup) = f(a) \tc \{\etup\} = f(a)$.
Note also that a morphism $\tL{\,} \to B$ of $\tL{\PreU}$ is just an upper set of $B$.
The poset $\U\tL{\,}$ has two elements, $\{\etup\}$ (``feasible'') and $\emptyset$ (``infeasible''), with $\{\etup\} \leq \emptyset$; so a morphism $f \colon A \to \tL{\,}$ is determined by the lower set $\{ a \mid f(a) = \{\etup\} \}$ of feasible functionalities.

\begin{remark}[Strict and non-strict]
    The inclusion $\tL{\PreU} \to \PreU$ is full and faithful, and every preorder $P$ is isomorphic in $\PreU$ to the tuple preorder $\tL{P}$, by $J_{\U}$ of the order isomorphism $p \mapsto \tup{p}$.
    So the inclusion is an equivalence of categories.
    Together with the isomorphisms $J_{\U}(\tup{a,b} \mapsto a \tc b) \colon P \times Q \to P \otimes Q$ of \cref{lem:splitting}, it relates the monoidal structure of \cref{thm:PreU-monoidal} to the strict one of \cref{thm:strict-monoidal}, in the same way as for $\tL{\Pre}$ and $\Pre$.
    We do not need this here.
\end{remark}

\begin{remark}[Is $\tL{\PreU}$ a Kleisli category?]
    The functor $\U$ does not map $\tL{\Pre}$ to itself, since the elements of $\U P$ are upper sets, not tuples.
    This is why \cref{def:PreUL} defines $\tL{\PreU}$ as a full subcategory of $\PreU$, in the same way as the book defines $\tL{\mRel}$ as a full subcategory of $\mRel$.
    If one prefers a Kleisli category, one can use the monad $P \mapsto \tL{\U P}$ on $\tL{\Pre}$, whose values are preorders of tuples of length one; its Kleisli category is isomorphic to $\tL{\PreU}$, by removing the brackets around each value.
\end{remark}

By \cref{lem:wf-product,lem:splitting}, a monoidal product of two well-founded tuple posets is well-founded.
So everything restricts to the full subcategory of $\tL{\PosU}$ on the well-founded posets, and, exactly as in \cref{prop:antichain-product}, in terms of antichains the monoidal product is
\begin{equation}
    \label{eq:strict-antichain-product}
    (h \otimes k)(a \tc c) = h(a) \tc k(c),
\end{equation}
and the braiding is $a \tc b \mapsto \{ b \tc a \}$.
We write $\tL{\PosA}$ for the resulting symmetric strict monoidal category: its objects are the well-founded tuple posets, and its morphisms, composition, and identities are those of $\PosA$ (\cref{prop:antichain-category}).

\section{Feedback}
\label{sec:feedback}

\subsection{Traces and traced functors}

Let $\tup{\Cat{C}, \otimes, I, \sigma}$ be a symmetric strict monoidal category.
Following \lbl{def:feedback-cat} and \lbl{def:trace-operator}, a \emph{trace} on $\Cat{C}$ is a family of functions
\begin{equation*}
    \Tr^{C}_{A,B} \colon \operatorname{Hom}_{\Cat{C}}(A \otimes C, B \otimes C) \to \operatorname{Hom}_{\Cat{C}}(A, B)
\end{equation*}
which satisfies the following conditions, for all morphisms for which the two sides are defined.
\begin{align}
    \label{eq:ax-nat-A}
    & \text{Naturality in $A$:} &
    \Tr^{C}_{A',B}\bigl((v \otimes \id_C) \then u\bigr) & = v \then \Tr^{C}_{A,B}(u)
    && (v \colon A' \to A) \\
    \label{eq:ax-nat-B}
    & \text{Naturality in $B$:} &
    \Tr^{C}_{A,B'}\bigl(u \then (v \otimes \id_C)\bigr) & = \Tr^{C}_{A,B}(u) \then v
    && (v \colon B \to B') \\
    \label{eq:ax-dinat}
    & \text{Dinaturality in $C$:} &
    \Tr^{C}_{A,B}\bigl(u \then (\id_B \otimes v)\bigr) & = \Tr^{C'}_{A,B}\bigl((\id_A \otimes v) \then u\bigr)
    && (v \colon C' \to C) \\
    \label{eq:ax-van-1}
    & \text{Vanishing I:} &
    \Tr^{I}_{A,B}(u) & = u \\
    \label{eq:ax-van-2}
    & \text{Vanishing II:} &
    \Tr^{C}_{A,B}\bigl(\Tr^{D}_{A \otimes C, B \otimes C}(u)\bigr) & = \Tr^{C \otimes D}_{A,B}(u) \\
    \label{eq:ax-sup}
    & \text{Superposing:} &
    \Tr^{C}_{E \otimes A, F \otimes B}(v \otimes u) & = v \otimes \Tr^{C}_{A,B}(u)
    && (v \colon E \to F) \\
    \label{eq:ax-yank}
    & \text{Yanking:} &
    \Tr^{C}_{C,C}(\sigma_{C,C}) & = \id_C
\end{align}
In \cref{eq:ax-nat-A,eq:ax-nat-B,eq:ax-sup} we have $u \colon A \otimes C \to B \otimes C$; in \cref{eq:ax-dinat}, $u \colon A \otimes C \to B \otimes C'$; in \cref{eq:ax-van-1}, $u \colon A \to B$; and in \cref{eq:ax-van-2}, $u \colon A \otimes C \otimes D \to B \otimes C \otimes D$.
(A \emph{feedback operator} satisfies \cref{eq:ax-nat-A,eq:ax-nat-B,eq:ax-van-1,eq:ax-van-2,eq:ax-sup}, and \cref{eq:ax-dinat} for isomorphisms $v$ only.)

\begin{definition}
    \label{def:traced-functor}
    Let $\Cat{C}$ and $\Cat{D}$ be symmetric strict monoidal categories, each equipped with a trace.
    A strict symmetric monoidal functor $G \colon \Cat{C} \to \Cat{D}$ is \emph{traced} if
    \begin{equation}
        \label{eq:traced-functor}
        G\bigl(\Tr^{C}_{A,B}(u)\bigr) = \Tr^{GC}_{GA,GB}(Gu)
    \end{equation}
    for all morphisms $u \colon A \otimes C \to B \otimes C$ of $\Cat{C}$.
    An \emph{isomorphism of traced symmetric strict monoidal categories} is a traced strict symmetric monoidal functor which is an isomorphism of categories.
\end{definition}

Note that the right-hand side of \cref{eq:traced-functor} makes sense because $G(A \otimes C) = GA \otimes GC$.

\subsection{The trace of monotone relations}

For objects $A$, $B$, $C$ of $\tL{\mRel}$ and a monotone relation $M \colon A \otimes C \profto B \otimes C$, the book defines (\lbl{eq:mRelL-feedback})
\begin{equation}
    \label{eq:mRelL-trace}
    \Tr^{C}_{A,B}(M) = \{\, \tup{a,b} \mid \exists\, c \in C \colon \rel{(a \tc c)}{M}{(b \tc c)} \,\} \colon A \profto B.
\end{equation}

\begin{proposition}
    \label{prop:mRelL-trace}
    \Cref{eq:mRelL-trace} defines a trace on the symmetric strict monoidal category $\tL{\mRel}$.
    Moreover, $\tup{a,b} \in \Tr^{C}_{A,B}(M)$ if and only if there are $c, c' \in C$ with $c' \leq c$ and $\rel{(a \tc c)}{M}{(b \tc c')}$.
\end{proposition}

The second statement is \lbl{eq:mRelL-feedback-leq}.
The book states the first one and checks the yanking condition; \cref{app:mrel-trace} contains a proof of all the conditions.

\subsection{The trace of queries}

\begin{theorem}
    \label{thm:trace}
    For objects $A$, $B$, $C$ of $\tL{\PreU}$ and a morphism $h \colon A \otimes C \to B \otimes C$ of $\tL{\PreU}$, that is, a monotone function $h \colon A \otimes C \to \U(B \otimes C)$, define
    \begin{equation}
        \label{eq:kleisli-trace}
        \Tr^{C}_{A,B}(h) \colon A \to \U B,
        \qquad
        \Tr^{C}_{A,B}(h)(a) = \{\, b \in B \mid \exists\, c \in C \colon\ b \tc c \in h(a \tc c) \,\}.
    \end{equation}
    \begin{enumerate}[label=(\alph*)]
        \item $\Tr^{C}_{A,B}(h)$ is a morphism $A \to B$ of $\tL{\PreU}$.
        \item These operators form a trace on the symmetric strict monoidal category $\tL{\PreU}$ of \cref{thm:strict-monoidal}.
        \item For every monotone relation $M \colon A \otimes C \profto B \otimes C$,
              \begin{equation}
                  \label{eq:hat-trace}
                  \widehat{\Tr^{C}_{A,B}(M)} = \Tr^{C}_{A,B}\bigl(\widehat{M}\bigr).
              \end{equation}
              That is, $\Phi \colon \tL{\mRel} \to \tL{\PreU}$ is an isomorphism of traced symmetric strict monoidal categories.
        \item The trace restricts to $\tL{\PosU}$, and $\Phi \colon \tL{\Feas} \to \tL{\PosU}$ is an isomorphism of traced symmetric strict monoidal categories.
    \end{enumerate}
\end{theorem}
\begin{proof}
    By \cref{lem:transport}(c), \cref{prop:mRelL-trace}, and \cref{thm:strict-monoidal}, the operators $\Tr'(h) = \Phi(\Tr(\Phi^{-1}h))$ form a trace on $\tL{\PreU}$, and it is the only one for which $\Phi$ is traced.
    We compute it.
    By \cref{eq:phi-inverse}, $\Phi^{-1}h$ relates $a \tc c$ to $b \tc c'$ precisely when $b \tc c' \in h(a \tc c)$.
    By \cref{eq:mRelL-trace}, $\tup{a,b} \in \Tr(\Phi^{-1}h)$ precisely when there is $c \in C$ with $b \tc c \in h(a \tc c)$.
    Hence
    \begin{equation*}
        \Tr'(h)(a) = \{\, b \mid \tup{a,b} \in \Tr(\Phi^{-1}h) \,\} = \{\, b \mid \exists\, c \colon b \tc c \in h(a \tc c) \,\},
    \end{equation*}
    which is \cref{eq:kleisli-trace}.
    This proves (a), (b), and (c).
    For (d), the trace does not change the objects involved, so it restricts to the full subcategories on the posets.
\end{proof}

\begin{corollary}[The query functor]
    \label{cor:fixfunminres}
    The functor
    \begin{equation*}
        \Fix \colon \tL{\DP} \to \tL{\PosU},
        \qquad
        d \mapsto \widehat{d},
        \qquad
        \widehat{d}(a) = \{\, b \mid d(a^*, b) = \top \,\},
    \end{equation*}
    is an isomorphism of traced symmetric strict monoidal categories.
    Explicitly, for design problems $d$ and $e$ for which the left-hand sides are defined,
    \begin{align}
        \label{eq:query-series}
        \widehat{d \then e}(a) & = \textstyle\bigcup_{b \in \widehat{d}(a)} \widehat{e}(b), &
        \widehat{\id_A}(a) & = \up a, \\
        \label{eq:query-parallel}
        \widehat{d \otimes e}(a \tc c) & = \widehat{d}(a) \tc \widehat{e}(c), &
        \widehat{\sigma_{A,B}}(a \tc b) & = \up (b \tc a), \\
        \label{eq:query-feedback}
        \widehat{\Tr^{C}_{A,B}(d)}(a) & = \{\, b \mid \exists\, c \colon b \tc c \in \widehat{d}(a \tc c) \,\}.
    \end{align}
\end{corollary}
\begin{proof}
    $\Fix$ is the composite of $\Phi \colon \tL{\Feas} \to \tL{\PosU}$ with the functor $\tL{\DP} \to \tL{\Feas}$ which sends a design problem to its feasible set.
    The latter is an isomorphism of categories (\lbl{def:DP}), it is compatible with the monoidal products (\lbl{eq:dpmulti}, \lbl{eq:RelL-monoidal-product}) and the braidings (\lbl{exa:RelL-symmetric}), and it is compatible with the traces (\lbl{eq:codesign-loop-feasible-set}).
    Now apply \cref{thm:trace}(d).
\end{proof}

So the answer to the query for a system which is built from components by series composition, parallel composition, and feedback can be computed from the answers for the components, by \cref{eq:query-series,eq:query-parallel,eq:query-feedback}.

\subsection{Reading the trace formula}

For a fixed functionality $a$, call $c \in C$ a \emph{loop value for $b$} if $b \tc c \in h(a \tc c)$: providing the functionality $c$ on the loop wire, in addition to $a$, is feasible with the requirements $b$ and with the \emph{same} $c$ as requirement on the loop wire.
Then \cref{eq:kleisli-trace} says that $b$ is a feasible requirement for $a$ in the system with the loop closed if and only if there is a loop value for $b$.
As for relations, one may relax the condition on the loop wire to an inequality.

\begin{proposition}
    \label{prop:trace-leq}
    Let $h \colon A \otimes C \to B \otimes C$ be a morphism of $\tL{\PreU}$, and let $a \in A$ and $b \in B$.
    The following are equivalent.
    \begin{enumerate}[label=(\roman*)]
        \item $b \in \Tr^{C}_{A,B}(h)(a)$.
        \item There are $c, c' \in C$ with $c' \leq c$ and $b \tc c' \in h(a \tc c)$.
    \end{enumerate}
    More precisely, if $c' \leq c$ and $b \tc c' \in h(a \tc c)$, then both $c$ and $c'$ are loop values for $b$.
\end{proposition}
\begin{proof}
    (i) $\Rightarrow$ (ii): take $c' = c$.
    For the last statement, which implies (ii) $\Rightarrow$ (i), suppose $c' \leq c$ and $b \tc c' \in h(a \tc c)$.
    Since $b \tc c' \leq b \tc c$ and $h(a \tc c)$ is an upper set, $b \tc c \in h(a \tc c)$.
    Since $a \tc c' \leq a \tc c$ and $h$ is monotone, $h(a \tc c) \subseteq h(a \tc c')$ by \cref{eq:kleisli-antitone}, so $b \tc c' \in h(a \tc c')$.
\end{proof}

The two halves of this argument are the two closure conditions of a monotone relation: one may \emph{require more} on the loop wire (closure upward in the target), or \emph{provide less} (closure downward in the source).

\begin{remark}[Direct proofs]
    \label{rem:direct-proofs}
    The conditions \cref{eq:ax-nat-A}--\cref{eq:ax-yank} for $\tL{\PreU}$ can of course also be checked directly from \cref{eq:kleisli-U}, \cref{eq:strict-kleisli-product}, and \cref{eq:kleisli-trace}.
    The draft chapter does this for $\Cat{Pos}_{\U}$ (\lbl{lem:UPos-is-traced}), with one step left open, namely naturality in $B$.
    With \cref{prop:trace-leq}, this step is as follows.
    For $h \colon A \otimes C \to \U(B \otimes C)$ and $g \colon B \to \U B'$, and $a \in A$,
    \begin{align*}
        \Tr\bigl(h \thenU (g \otimes \id_C)\bigr)(a)
        & = \{\, b' \mid \exists\, c,\ b,\ c' \colon\ b \tc c' \in h(a \tc c),\ b' \in g(b),\ c' \leq c \,\}, \\
        \bigl(\Tr(h) \thenU g\bigr)(a)
        & = \{\, b' \mid \exists\, c,\ b \colon\ b \tc c \in h(a \tc c),\ b' \in g(b) \,\}.
    \end{align*}
    The second set is included in the first (take $c' = c$), and the first is included in the second by \cref{prop:trace-leq}.
    The point is that the identity morphism of $C$ in $\tL{\PreU}$ is $c' \mapsto \up c'$ and not $c' \mapsto \{c'\}$, which produces the inequality $c' \leq c$, and that this inequality can be absorbed because $h$ is monotone with values in upper sets.
\end{remark}

\subsection{Feedback in terms of antichains}

\begin{proposition}
    \label{prop:trace-antichains}
    Let $A$, $B$, and $C$ be well-founded tuple posets, and let $h \colon A \otimes C \to \An(B \otimes C)$ be a morphism of $\tL{\PosA}$, with corresponding morphism $f = h \then \up$ of $\tL{\PosU}$.
    Then, for every $a \in A$,
    \begin{equation}
        \label{eq:antichain-trace}
        \Min \Tr^{C}_{A,B}(f)(a) = \Min \{\, b \in B \mid \exists\, c \in C \colon\ b \tc c \in h(a \tc c) \,\}.
    \end{equation}
    So, in terms of antichains, the trace is given by the same formula as \cref{eq:kleisli-trace}, followed by $\Min$.
\end{proposition}
\begin{proof}
    Let $X_a = \{ b \mid \exists\, c \colon b \tc c \in h(a \tc c) \}$.
    We show $\Tr(f)(a) = \up X_a$; then \cref{eq:antichain-trace} follows from \cref{lem:min-up}(b).

    ``$\supseteq$'': if $b_0 \tc c \in h(a \tc c)$, then $b_0 \tc c \in f(a \tc c)$, so $b_0 \in \Tr(f)(a)$; and $\Tr(f)(a)$ is an upper set.

    ``$\subseteq$'': let $b \in \Tr(f)(a)$.
    The set $K = \{ c \in C \mid b \tc c \in f(a \tc c) \}$ is non-empty, so, since $C$ is well-founded, it has a minimal element $c$.
    Since $b \tc c \in f(a \tc c) = \up h(a \tc c)$, there is an element $b_0 \tc c_0 \in h(a \tc c)$ with $b_0 \leq b$ and $c_0 \leq c$ (\cref{lem:splitting}).
    Then $b \tc c_0 \in f(a \tc c)$, because $f(a \tc c)$ is an upper set, and so $c_0 \in K$ by \cref{prop:trace-leq}.
    Since $c$ is minimal in $K$ and $c_0 \leq c$, we get $c_0 = c$.
    Hence $b_0 \tc c \in h(a \tc c)$, that is, $b_0 \in X_a$, and $b_0 \leq b$.
\end{proof}

In words: to close a loop on the level of Pareto fronts, go through the loop values $c$; in the Pareto front $h(a \tc c)$ of the open system, keep the points whose requirement on the loop wire is \emph{exactly} $c$; forget the loop coordinate; and finally discard the dominated points.
Points $b \tc c'$ of $h(a \tc c)$ with $c' < c$ need not be considered separately: by \cref{prop:trace-leq}, such a $b$ lies in $\Tr(f)(a) = \up X_a$, so it is above a point which is found in the way just described.
\Cref{exa:two-modules} shows that the final $\Min$ is needed.

\subsection{The dual query}

The isomorphism $\Psi$ of \cref{thm:mrel-kleisli-lower} is compatible with parallel composition and feedback in the same way.

\begin{proposition}
    \label{prop:dual-monoidal-trace}
    For monotone relations $M \colon A \profto B$ and $N \colon C \profto D$ in $\tL{\mRel}$,
    \begin{equation*}
        \widecheck{M \otimes N}(b \tc d) = \widecheck{M}(b) \tc \widecheck{N}(d),
    \end{equation*}
    and for a monotone relation $M \colon A \otimes C \profto B \otimes C$,
    \begin{equation*}
        \widecheck{\Tr^{C}_{A,B}(M)}(b) = \{\, a \in A \mid \exists\, c \in C \colon\ a \tc c \in \widecheck{M}(b \tc c) \,\}.
    \end{equation*}
\end{proposition}
\begin{proof}
    By \cref{eq:mRelL-product}, $a \tc c \in \widecheck{M \otimes N}(b \tc d)$ if and only if $\rel{a}{M}{b}$ and $\rel{c}{N}{d}$.
    By \cref{eq:mRelL-trace}, $a \in \widecheck{\Tr(M)}(b)$ if and only if there is $c$ with $\rel{(a \tc c)}{M}{(b \tc c)}$, that is, with $a \tc c \in \widecheck{M}(b \tc c)$.
\end{proof}

So the right-hand sides define a monoidal product and a trace on the full subcategory $\tL{\PreL}$ of $\PreL$ on the tuple preorders, and $\Psi$ identifies $\tL{\mRel}\op$, with the same monoidal product and trace operators as $\tL{\mRel}$, with $\tL{\PreL}$.
If the posets involved are \emph{dually} well-founded (every non-empty subset has a maximal element), lower sets correspond to antichains via $\dn$ and $\Max$, and the statements of \cref{prop:antichain-category,prop:antichain-product,prop:trace-antichains} hold with $\Max$ in place of $\Min$.

\begin{remark}[Computation]
    \Cref{eq:kleisli-trace} is a specification, not an algorithm: it quantifies over all elements of $C$.
    If $C$ is finite, it is a finite search, and \cref{prop:trace-antichains} reduces it to a search for fixed points on Pareto fronts.
    For infinite posets, such as the real numbers, the minimal requirements of a system with a loop are computed in \cite{censi} as a least fixed point, by an iteration which converges under continuity assumptions; see also the draft section \file{60_computation/20_solving/18_handling_loops.tex}.
    The relation of that iteration to \cref{prop:trace-antichains} is not treated in this note.
\end{remark}

\section{Worked examples}
\label{sec:examples}

In the examples we write tuples of length one as numbers.

\begin{example}[Series: paying with coins, and carrying them]
    \label{exa:coins}
    Let $\N$ be the natural numbers with their usual order.
    As in \lbl{exa:coins-dp}, consider the feasibility relation
    \begin{equation*}
        \mathit{pay} \colon \tL{\N} \profto \tL{\N, \N},
        \qquad
        \rel{a}{\mathit{pay}}{\tup{x,y}} \iff a \leq x + 2y,
    \end{equation*}
    which says that $x$ one-franc coins and $y$ two-franc coins suffice to pay the amount $a$.
    Suppose that a one-franc coin weighs $5$ grams and a two-franc coin weighs $8$ grams, and consider
    \begin{equation*}
        \mathit{carry} \colon \tL{\N, \N} \profto \tL{\N},
        \qquad
        \rel{\tup{x,y}}{\mathit{carry}}{w} \iff 5x + 8y \leq w,
    \end{equation*}
    which says that a purse for $w$ grams can carry the coins.
    The series composition $\mathit{pay} \then \mathit{carry} \colon \tL{\N} \profto \tL{\N}$ relates an amount to the purses with which it can be paid.

    All posets involved are well partial orders, so by \cref{prop:wpo} all queries have finite antichains of minimal answers.
    For the first component, the set $\widehat{\mathit{pay}}(a) = \{ \tup{x,y} \mid a \leq x + 2y \}$ has the minimal elements
    \begin{equation*}
        h_{\mathit{pay}}(a) = \{\, \tup{a - 2j,\, j} \mid 0 \leq j \leq \lfloor a/2 \rfloor \,\} \cup \{\, \tup{0, (a+1)/2} \mid a \text{ odd} \,\};
    \end{equation*}
    for instance $h_{\mathit{pay}}(3) = \{ \tup{3,0}, \tup{1,1}, \tup{0,2} \}$, in agreement with \lbl{exa:coins-dp}.
    For the second component, $\widehat{\mathit{carry}}\tup{x,y} = \up (5x + 8y)$ has the least element $5x + 8y$, so $h_{\mathit{carry}}\tup{x,y} = \{ 5x + 8y \}$; this is a companion (\cref{prop:companion-free}).
    By \cref{eq:antichain-composition},
    \begin{equation*}
        (h_{\mathit{pay}} \thenA h_{\mathit{carry}})(3) = \Min\bigl( \{15\} \cup \{13\} \cup \{16\} \bigr) = \{13\}.
    \end{equation*}
    The three ways of paying three francs give the weights $15$, $13$, and $16$, and $\Min$ discards the dominated ones: the lightest way to pay three francs is one coin of each kind.
    In the same way one finds, for every $a$,
    \begin{equation*}
        (h_{\mathit{pay}} \thenA h_{\mathit{carry}})(a) = \begin{cases} \{4a\} & \text{if $a$ is even,} \\ \{4a + 1\} & \text{if $a$ is odd.} \end{cases}
    \end{equation*}
    In terms of upper sets, \cref{eq:kleisli-U} gives $\widehat{\mathit{pay} \then \mathit{carry}}(3) = \bigcup_{3 \leq x + 2y} \up (5x + 8y) = \up 13$, a union of infinitely many upper sets; the antichain form needs only the three minimal elements.
\end{example}

\begin{example}[Parallel]
    \label{exa:coins-parallel}
    To pay $3$ francs to one person and $2$ francs to another, we use $\mathit{pay} \otimes \mathit{pay} \colon \tL{\N,\N} \profto \tL{\N,\N,\N,\N}$.
    Since $h_{\mathit{pay}}(2) = \{ \tup{2,0}, \tup{0,1} \}$, \cref{eq:strict-antichain-product} gives
    \begin{align*}
        (h_{\mathit{pay}} \otimes h_{\mathit{pay}})(3 \tc 2)
        & = h_{\mathit{pay}}(3) \tc h_{\mathit{pay}}(2) \\
        & = \{ \tup{3,0,2,0}, \tup{3,0,0,1}, \tup{1,1,2,0}, \tup{1,1,0,1}, \tup{0,2,2,0}, \tup{0,2,0,1} \},
    \end{align*}
    an antichain with $3 \cdot 2 = 6$ elements; no element is discarded.
    The functoriality of the monoidal product \cref{eq:kleisli-interchange} can be seen at work here: computing $(\mathit{pay} \then \mathit{carry}) \otimes (\mathit{pay} \then \mathit{carry})$ at $3 \tc 2$ gives $\{13\} \tc \{8\} = \{ \tup{13,8} \}$, and computing $(\mathit{pay} \otimes \mathit{pay}) \then (\mathit{carry} \otimes \mathit{carry})$ gives the minimal elements of the six pairs of weights $\tup{15,10}, \tup{15,8}, \tup{13,10}, \tup{13,8}, \tup{16,10}, \tup{16,8}$, which is again $\{ \tup{13,8} \}$.
\end{example}

\begin{example}[Feedback: battery modules]
    \label{exa:modules}
    We redo \lbl{exa:codesign-feedback-modules} in terms of queries.
    Let $A = \tL{\N}$, $C = \tL{\{0,1,2,3\}}$, and $B = \tL{\,}$, so that $B \otimes C = C$, and let
    \begin{equation*}
        M \colon A \otimes C \profto C,
        \qquad
        \rel{(a \tc c)}{M}{c'} \iff a + c \leq 2c'.
    \end{equation*}
    Here $a$ is the mass of the payload, $c$ is the number of battery modules on board, and $c'$ is the number of modules which are needed to move both; each module has mass one and can move two units of mass, and there is room for at most three modules.
    The query of the open system is
    \begin{align*}
        \widehat{M}(a \tc c) & = \{\, c' \in C \mid a + c \leq 2c' \,\}, \\
        h(a \tc c) = \Min \widehat{M}(a \tc c) & =
        \begin{cases}
            \{ \lceil (a+c)/2 \rceil \} & \text{if } \lceil (a+c)/2 \rceil \leq 3, \\
            \emptyset & \text{otherwise.}
        \end{cases}
    \end{align*}
    By \cref{eq:kleisli-trace}, $\Tr^{C}_{A,B}(\widehat{M})(a)$ is $\{\etup\}$ if there is a loop value $c$, that is, a $c \leq 3$ with $c \in \widehat{M}(a \tc c)$, and $\emptyset$ otherwise.
    The condition $c \in \widehat{M}(a \tc c)$ reads $a + c \leq 2c$, that is, $a \leq c$.
    Hence
    \begin{equation*}
        \Tr^{C}_{A,B}(\widehat{M})(a) =
        \begin{cases}
            \{\etup\} & \text{if } a \leq 3, \\
            \emptyset & \text{if } a > 3,
        \end{cases}
    \end{equation*}
    which is the answer found in the book: the vehicle with the loop closed can transport any payload of mass at most three, and has no requirements.

    In terms of antichains (\cref{prop:trace-antichains}), we look for $c$ with $c \in h(a \tc c)$, that is, $c = \lceil (a+c)/2 \rceil$, which holds precisely for $c = a$ and $c = a + 1$.
    So for $a = 1$, the loop values found on the Pareto fronts are $c = 1$ and $c = 2$, while in the sense of upper sets all of $c = 1, 2, 3$ are loop values.
\end{example}

\begin{example}[Feedback with a Pareto front: two kinds of modules]
    \label{exa:two-modules}
    Suppose now that there are two kinds of battery modules.
    Both have mass one and can move two units of mass, and there is room for at most three modules of each kind.
    A module of the first kind costs $1$ and causes emissions $3$; a module of the second kind costs $2$ and causes emissions $1$.
    Let
    \begin{equation*}
        A = \tL{\N},
        \qquad
        C = \tL{\{0,1,2,3\}, \{0,1,2,3\}},
        \qquad
        B = \tL{\N, \N},
    \end{equation*}
    where the elements of $B$ are pairs of a price and an amount of emissions, and let $M \colon A \otimes C \profto B \otimes C$ relate $a \tc \tup{c_1, c_2}$ to $\tup{p, e} \tc \tup{c_1', c_2'}$ if and only if
    \begin{equation*}
        a + c_1 + c_2 \leq 2(c_1' + c_2'),
        \qquad
        c_1 + 2c_2 \leq p,
        \qquad
        3c_1 + c_2 \leq e.
    \end{equation*}
    As in \lbl{exa:codesign-feedback-vehicle}, the price and the emissions are those of the modules on board, and the modules which are needed are a requirement; one checks that $M$ is closed downward in the source and upward in the target.
    The Pareto fronts of the open system are
    \begin{equation*}
        h(a \tc \tup{c_1, c_2}) = \bigl\{\, \tup{c_1 + 2c_2,\ 3c_1 + c_2} \tc \tup{c_1', c_2'} \bigm| c_1' + c_2' = \lceil (a + c_1 + c_2)/2 \rceil,\ c_1', c_2' \leq 3 \,\bigr\};
    \end{equation*}
    for instance $h(2 \tc \tup{1,1}) = \{ \tup{3,4,0,2}, \tup{3,4,1,1}, \tup{3,4,2,0} \}$.
    A point of $h(a \tc c)$ has the loop requirement $c$ precisely when $c_1 + c_2 = \lceil (a + c_1 + c_2)/2 \rceil$, that is, when $c_1 + c_2 \in \{a, a+1\}$.
    So the set $X_a$ of \cref{prop:trace-antichains} is
    \begin{equation*}
        X_a = \{\, \tup{c_1 + 2c_2,\ 3c_1 + c_2} \mid c_1, c_2 \leq 3,\ c_1 + c_2 \in \{a, a+1\} \,\}.
    \end{equation*}
    For $a = 2$ this is
    \begin{equation*}
        X_2 = \{ \tup{2,6}, \tup{3,4}, \tup{4,2} \} \cup \{ \tup{3,9}, \tup{4,7}, \tup{5,5}, \tup{6,3} \},
    \end{equation*}
    where the first three points come from two modules on board and the last four from three modules.
    Each of the last four is dominated by one of the first three, so
    \begin{equation*}
        \Min \Tr^{C}_{A,B}(\widehat{M})(2) = \Min X_2 = \{ \tup{2,6}, \tup{3,4}, \tup{4,2} \}.
    \end{equation*}
    This shows that the $\Min$ in \cref{eq:antichain-trace} cannot be omitted.
    In general, the minimal requirements of the vehicle with the loop closed are obtained from exactly $a$ modules, and they form a trade-off between price and emissions:
    \begin{center}
        \renewcommand{\arraystretch}{1.2}
        \begin{tabular}{cl}
            \toprule
            payload $a$ & minimal requirements $\tup{\text{price}, \text{emissions}}$ \\
            \midrule
            $0$ & $\tup{0,0}$ \\
            $1$ & $\tup{1,3}$, $\tup{2,1}$ \\
            $2$ & $\tup{2,6}$, $\tup{3,4}$, $\tup{4,2}$ \\
            $3$ & $\tup{3,9}$, $\tup{4,7}$, $\tup{5,5}$, $\tup{6,3}$ \\
            $4$ & $\tup{5,10}$, $\tup{6,8}$, $\tup{7,6}$ \\
            $5$ & $\tup{7,11}$, $\tup{8,9}$ \\
            $6$ & $\tup{9,12}$ \\
            $\geq 7$ & none (infeasible) \\
            \bottomrule
        \end{tabular}
    \end{center}
    As the payload grows, the antichains move up in the order \cref{eq:antichain-order}, until they reach the empty antichain.
\end{example}

\section{Integrating the material into the book}
\label{sec:integration}

The story has four layers, which fit four different places in the book: formulas for composing answers to queries, which need only order theory; the statement that passing to queries is a functor, and an isomorphism of categories; its compatibility with parallel composition and feedback; and the description in terms of a monad.
The following suggestions go through the book in order.
File names are relative to \file{volumes/vol1}.

\paragraph{Part ``Order'', chapter on monotone relations (\file{05_order/50_monotone_relations}).}
The function $\widehat{M}$ is already introduced there, before \lbl{lem:which-monrels-are-companions}.
Two short lemmas would prepare everything that follows.
First, in the subsection on composing monotone relations, next to the matrix formula \lbl{lem:boolprof-composition-matrix}: the formula $\widehat{M \then N}(p) = \bigcup_{q \in \widehat{M}(p)} \widehat{N}(q)$ and $\widehat{\leq_P}(p) = \up p$ (the last two steps of the proof of \cref{thm:mrel-kleisli}).
Second, in \file{40_products_sums.tex}, after \lbl{lem:monrel-product}: the formula $\widehat{M \times N}\tup{p,q} = \widehat{M}(p) \times \widehat{N}(q)$, together with \cref{lem:upper-closure}(e).
The observation after \cref{thm:mrel-kleisli}, that the two closure conditions of a monotone relation are exactly ``values are upper sets'' and ``the function is monotone'', would fit the summary at the end of \file{10_monotone_relations.tex}.

\paragraph{Part ``Order'', chapter on design problems (\file{05_order/60_design_problems}).}
In \file{40_optimality.tex}, subsection ``Antichains and upper sets'':
\begin{itemize}
    \item \lbl{lem:finite-upper-sets-dc} can be stated for well-founded posets instead of finite ones, with a shorter proof (the first half of the proof of \cref{thm:antichains-wf}), and \cref{thm:antichains-wf} answers the question of when exactly upper sets and antichains correspond.
          This also covers the running example: for paying with coins (\lbl{exa:coins-dp}) the poset of requirements $\N \times \N$ is infinite, so the present lemma does not apply to it, although its sets of feasible requirements are generated by the minimal elements which the text computes; with \cref{lem:wf-product} it is covered.
          The sensor catalogue (\lbl{exa:sensor-catalogue}) has requirements in $\Rp \times \Rp$, which is not well-founded; there the sets of feasible requirements are upper closures of finite sets by construction, so \lbl{prop:antichains-upper-sets} applies directly.
    \item \cref{lem:min-up}(b), $\Min \up X = \Min X$, is worth stating; with it, the last formula of \lbl{prop:dp-series-queries} simplifies to $\Min \bigcup_{q} h_e(q)$ (\cref{rem:book-series-queries}).
    \item A subsection ``Answering queries for parallel compositions'' could follow the one for series compositions, with \cref{lem:min-product} and \cref{eq:antichain-product}, once products of design problems are available in that chapter (at present only the product of monotone relations is defined in the Order part).
    \item \Cref{prop:wpo} and \cref{exa:closed-upper-sets} could be remarks: the first explains why answers are finite for requirements in $\N^k$, the second what replaces well-foundedness for real-valued requirements.
    \item A remark on terminology: the name ``downward-closed upper set'' of \lbl{def:downward-closed-upperset} may be read as ``upper set which is also a lower set''. A phrase such as ``upper set generated by its minimal elements'' says what is meant.
\end{itemize}
\Cref{exa:coins} could serve as a second example for \lbl{prop:dp-series-queries}, since it continues the running example of that chapter.

\paragraph{Parts ``Categories'' and ``Functors''.}
The section on codesign queries in the chapter of extended examples of that part,
\begin{center}
    \file{25_functors/30_application_examples_functors/15_codesign-querying.tex}
\end{center}
(\lbl{sec:functors-at-work-codesign-queries}), currently consists of a to-do item asking for the functors from design problems to the two query types.
\Cref{thm:mrel-kleisli,thm:mrel-kleisli-lower} are these functors, and they can be presented there without monads: define the category of queries directly by the formulas of \cref{prop:kleisli-U} (as the draft \lbl{def:upos_cat} does), and state that $\Fix \colon \DP \to \Cat{Pos}_{\U}$ is a functor and an isomorphism of categories.
The message for the reader is that functoriality means precisely that the answer for a series composition can be computed from the answers for the components.
$\FixR$ is a natural example of a \emph{contravariant} functor, $\DP\op \to \Cat{Pos}_{\Lo}$, and \cref{prop:companion-free} is an example of a commuting triangle of functors.
The category of queries could also be mentioned as a third description in the remark after \lbl{def:DP} in \file{10_categories/15_modeling_with_cats/46_series.tex}.
Possible graded exercises: the category axioms for $\Cat{Pos}_{\U}$ (unit laws use both ``upper set'' and ``monotone''); functoriality of $d \mapsto \widehat{d}$; the contravariant version.

\paragraph{Part ``Interconnection'', parallel composition (\file{27_interconnections/20_parallelism}).}
The category $\tL{\Cat{Pos}_{\U}}$ with the product \cref{eq:strict-kleisli-product} can be added to the paragraph on $\tL{\Rel}$, $\tL{\mRel}$, $\tL{\BoolProf}$ in \file{24_monoidal.tex}, or to the example on codesign (\lbl{exa:codesign-strict-monoidal}): the queries of two components side by side.
The section \file{50_symmetric_monoidal_functors.tex} is a stub.
The definition \cref{eq:strict-monoidal-functor} and three examples from this note would fill it: the companion functor $\tL{\Pre} \to \tL{\mRel}$, the passage $\tL{\DP} \to \tL{\Feas}$ from design problems to feasible sets, and $\Fix \colon \tL{\Feas} \to \tL{\Cat{Pos}_{\U}}$ (\cref{thm:strict-monoidal}).
\Cref{lem:transport} is a useful general tool there: it shows that $\tL{\Cat{Pos}_{\U}}$ is symmetric strict monoidal without checking any axiom.
\Cref{app:mrel-strict} is a proof of the statement about $\tL{\mRel}$ which the book currently makes without proof.

\paragraph{Part ``Interconnection'', feedback (\file{27_interconnections/50_feedback}).}
\begin{itemize}
    \item In \file{10_feedback_and_trace.tex}, paragraph ``A trace for monotone relations'': the sentence that the conditions ``can be checked in a similar way as for $\tL{\Rel}$'' can be backed by \cref{app:mrel-trace}, either as a proof or as graded exercises. Naturality in $B$ and dinaturality are instructive, because they are the places where the closure conditions of a monotone relation are used; for plain relations they are not needed.
    \item In the paragraph on codesign (\lbl{sec:feedback-codesign}): the query form \cref{eq:query-feedback} of closing the loop, the notion of a loop value, and \cref{prop:trace-leq}. \Cref{exa:modules} is the example \lbl{exa:codesign-feedback-modules} of the book in query form, and \cref{exa:two-modules} is a variant in which the answer is a Pareto front.
    \item The section \file{45_traced_monoidal_functors.tex} is a stub. \Cref{def:traced-functor} and \cref{cor:fixfunminres} would be its definition and its main example: closing a loop commutes with passing to queries.
    \item The draft section \file{20_closing_loop_codesign.tex} contains a longer verification of the trace axioms for design problems in the non-strict setting; \cref{app:mrel-trace} could replace it.
\end{itemize}

\paragraph{Part on monads (\file{34_monads-comonads/10_monads}).}
In \file{20_monads.tex}, after the example of the powerset monad: the upper set monad (\cref{def:U}, \cref{thm:U-monad}) and the lower set monad as examples, with \cref{rem:order} on the order.
In \file{30_kleisli.tex}, after the graded exercise \lbl{ex:HwkRelKleisli}: \cref{thm:mrel-kleisli} as its ordered analogue, and \cref{rem:discrete}.
Note that \cref{sec:warm-up} of this note is a solution of that graded exercise, so it should not go into the text.
\Cref{rem:monoidal-monad} is optional there.

\paragraph{Part ``Compositional computation'' (\file{60_computation}, draft).}
The chapters \file{15_query_equiv} and \file{40_uncertainty/30_LU_monads.tex} already contain $\Cat{Pos}_{\U}$, $\Cat{Pos}_{\Lo}$, a monoidal product, a trace, and the monad $\U$.
This note can be used to complete and shorten them.
\begin{itemize}
    \item \lbl{lem:upos_moncat} (proof missing), \lbl{lem:UPos-is-sym-mon} (a long direct check of the hexagon), and \lbl{lem:UPos-is-traced} (naturality in $B$ left open) all follow from the isomorphism with monotone relations by \cref{lem:transport}. If the direct proof of the trace axioms is kept, \cref{rem:direct-proofs} supplies the missing step.
    \item The lemma which states that $\Fix$ and its inverse form an \emph{equivalence}, with a to-do item about naturality: the two functors are inverse to each other on the nose, so this is an isomorphism of categories and no natural isomorphisms are needed.
    \item \lbl{lem:confunctor} should say that $\FixR$ is a functor $\DP\op \to \Cat{Pos}_{\Lo}$ (\cref{rem:U-L}).
    \item In \file{15_query_equiv/05_intro.tex}, the function $h_d$ is typed as a monotone function $F\op \to \tup{\An R, \subseteq}$. Inclusion is not the right order on antichains: in \lbl{exa:sensor-catalogue}, $h(30) = \{\tup{50,1}\}$ and $h(50) = \{\tup{100,5}, \tup{200,2}\}$ are not related by inclusion. It should be $F \to \An R$ with the order of \lbl{def:antichain-order}.
    \item In \file{30_LU_monads.tex}, the order on $\U P$ is not specified (it must be reverse inclusion, \cref{rem:order}), and the proof that $\U$ is a monad has open to-do items; the proof of \cref{thm:U-monad} can replace it. \lbl{lem:unpack_u_functor} is \cref{lem:upper-closure}(d), in the form $\up f[\up X] = \up f[X]$ suggested in the to-do item there.
    \item \lbl{lem:uposkleisli} (``$\Cat{Pos}_{\U}$ is the Kleisli category of $\U$'') is \cref{prop:kleisli-U}; the notation $\Cat{Pos}_{\U}$ already agrees with the notation $\Cat{C}_T$ of the book for Kleisli categories.
\end{itemize}

\paragraph{What this note does not cover.}
The relation between the trace formula and the fixed point iteration of \cite{censi}; a trace for the non-strict monoidal category $\PreU$ (the formula is the same, with pairs instead of concatenated tuples, and the non-strict axioms would be obtained by the same transport argument from $\mRel$); and the monoidal equivalence between $\tL{\PreU}$ and $\PreU$.

\appendix

\section{The symmetric strict monoidal category of monotone relations}
\label{app:mrel-strict}

We prove \cref{prop:mRelL}.
Throughout, $A, B, C, \ldots$ are objects of $\tL{\mRel}$, that is, tuple preorders, and we use \cref{lem:splitting} without mention: every element of $A \otimes C$ is of the form $a \tc c$ for unique $a \in A$ and $c \in C$, and tuples are compared entry by entry.

\paragraph{The product is well defined.}
Let $M \colon A \profto B$ and $N \colon C \profto D$ be monotone relations.
If $a' \tc c' \leq a \tc c$ and $\rel{(a \tc c)}{(M \otimes N)}{(b \tc d)}$, then $a' \leq a$, $\rel{a}{M}{b}$, $c' \leq c$, and $\rel{c}{N}{d}$, so $\rel{a'}{M}{b}$ and $\rel{c'}{N}{d}$ by \cref{eq:closed-source}, that is, $\rel{(a' \tc c')}{(M \otimes N)}{(b \tc d)}$.
Closure upward in the target is proved in the same way from \cref{eq:closed-target}.

\paragraph{The product is a functor.}
It preserves identities: ${\leq_A} \otimes {\leq_C}$ relates $a \tc c$ to $a' \tc c'$ precisely when $a \leq a'$ and $c \leq c'$, that is, when $a \tc c \leq a' \tc c'$.
It preserves composition: for $M \colon A \profto A'$, $M' \colon A' \profto A''$, $N \colon C \profto C'$, and $N' \colon C' \profto C''$, both $(M \then M') \otimes (N \then N')$ and $(M \otimes N) \then (M' \otimes N')$ relate $a \tc c$ to $a'' \tc c''$ precisely when there are $a' \in A'$ and $c' \in C'$ with $\rel{a}{M}{a'}$, $\rel{a'}{M'}{a''}$, $\rel{c}{N}{c'}$, and $\rel{c'}{N'}{c''}$.

\paragraph{Strictness.}
On objects, the monoidal product is that of $\tL{\Pre}$, which is associative and has $\tL{\,}$ as a neutral element.
For monotone relations $M$, $N$, $L$, both $(M \otimes N) \otimes L$ and $M \otimes (N \otimes L)$ relate $a \tc b \tc c$ to $a' \tc b' \tc c'$ precisely when $\rel{a}{M}{a'}$, $\rel{b}{N}{b'}$, and $\rel{c}{L}{c'}$, because concatenation is associative.
The identity of $\tL{\,}$ is $\{\tup{\etup, \etup}\}$, so $M \otimes \id_{\tL{\,}}$ relates $a \tc \etup = a$ to $b \tc \etup = b$ precisely when $\rel{a}{M}{b}$; hence $M \otimes \id_{\tL{\,}} = M$, and likewise $\id_{\tL{\,}} \otimes M = M$.

\paragraph{Companions.}
For monotone functions $u \colon A \to B$ and $v \colon C \to D$ between tuple preorders, we have
\begin{equation}
    \label{eq:companion-strict-product}
    \comp{(u \otimes v)} = \comp{u} \otimes \comp{v},
\end{equation}
since both sides relate $a \tc c$ to $b \tc d$ precisely when $u(a) \leq b$ and $v(c) \leq d$.
Together with $\comp{(\id)} = {\leq}$ and $\comp{(u \then v)} = \comp{u} \then \comp{v}$, this says that every equation between monotone functions which is built from composition, identities, and $\otimes$ gives an equation between their companions.

\paragraph{The braiding.}
Let $s_{A,B} \colon A \otimes B \to B \otimes A$ be the order isomorphism $a \tc b \mapsto b \tc a$.
The relation \cref{eq:mRelL-braiding} is its companion: $\sigma_{A,B}$ relates $a \tc b$ to $b' \tc a'$ precisely when $b \tc a \leq b' \tc a'$.
In particular it is a monotone relation.
The functions $s$ satisfy
\begin{equation*}
    s_{A,B} \then s_{B,A} = \id_{A \otimes B},
    \qquad
    s_{A, B \otimes C} = (s_{A,B} \otimes \id_C) \then (\id_B \otimes s_{A,C}),
\end{equation*}
as one sees by evaluating both sides at $a \tc b$ and at $a \tc b \tc c$.
Taking companions, we obtain the symmetry condition $\sigma_{A,B} \then \sigma_{B,A} = \id_{A \otimes B}$, which also shows that $\sigma_{A,B}$ is an isomorphism, and the compatibility with nesting $\sigma_{A, B \otimes C} = (\sigma_{A,B} \otimes \id_C) \then (\id_B \otimes \sigma_{A,C})$.

It remains to prove naturality.
Let $M \colon A \profto A'$ and $N \colon B \profto B'$ be monotone relations.
By \cref{lem:sliding}, both
\begin{equation*}
    (M \otimes N) \then \sigma_{A',B'}
    \qquad \text{and} \qquad
    \sigma_{A,B} \then (N \otimes M)
\end{equation*}
relate $a \tc b$ to $b' \tc a'$ precisely when $\rel{a}{M}{a'}$ and $\rel{b}{N}{b'}$: for the first relation, by \cref{eq:sliding-right}, the condition is $\rel{(a \tc b)}{(M \otimes N)}{s_{A',B'}^{-1}(b' \tc a')}$, and for the second, by \cref{eq:sliding-left}, it is $\rel{s_{A,B}(a \tc b)}{(N \otimes M)}{(b' \tc a')}$.

\paragraph{Restriction to posets.}
A monoidal product of two tuple posets is a poset, the unit is a poset, and the braidings between tuple posets are morphisms of $\tL{\Feas}$; so the structure restricts to $\tL{\Feas}$.
This completes the proof of \cref{prop:mRelL}. \qed

\section{The trace axioms for monotone relations}
\label{app:mrel-trace}

We prove \cref{prop:mRelL-trace}.
The conventions are those of \cref{app:mrel-strict}.
Recall that the identity morphism of $C$ is $\leq_C$, so that for a monotone relation $N \colon A' \profto A$,
\begin{equation}
    \label{eq:tensor-with-identity}
    \rel{(a' \tc c)}{(N \otimes \id_C)}{(a \tc c_1)} \iff \rel{a'}{N}{a} \text{ and } c \leq c_1.
\end{equation}

\paragraph{The trace is well defined, and the second description.}
Let $M \colon A \otimes C \profto B \otimes C$ be a monotone relation.
If $\tup{a,b} \in \Tr^{C}_{A,B}(M)$, with $\rel{(a \tc c)}{M}{(b \tc c)}$, and $a' \leq a$, then $a' \tc c \leq a \tc c$, so $\rel{(a' \tc c)}{M}{(b \tc c)}$ by \cref{eq:closed-source}, and $\tup{a',b} \in \Tr^{C}_{A,B}(M)$.
In the same way, $b \leq b'$ implies $\tup{a,b'} \in \Tr^{C}_{A,B}(M)$.
So $\Tr^{C}_{A,B}(M)$ is a monotone relation.
If $c' \leq c$ and $\rel{(a \tc c)}{M}{(b \tc c')}$, then $b \tc c' \leq b \tc c$, so $\rel{(a \tc c)}{M}{(b \tc c)}$ by \cref{eq:closed-target}; this proves the second statement of \cref{prop:mRelL-trace}.

\paragraph{Naturality in $A$.}
Let $M \colon A \otimes C \profto B \otimes C$ and $N \colon A' \profto A$.
By \cref{eq:tensor-with-identity}, $\tup{a', b}$ is an element of $\Tr^{C}_{A',B}((N \otimes \id_C) \then M)$ if and only if
\begin{equation}
    \label{eq:nat-A-lhs}
    \text{there are } c, c_1 \in C \text{ and } a \in A \text{ with } \rel{a'}{N}{a},\quad c \leq c_1, \quad \rel{(a \tc c_1)}{M}{(b \tc c)},
\end{equation}
and it is an element of $N \then \Tr^{C}_{A,B}(M)$ if and only if
\begin{equation}
    \label{eq:nat-A-rhs}
    \text{there are } c \in C \text{ and } a \in A \text{ with } \rel{a'}{N}{a}, \quad \rel{(a \tc c)}{M}{(b \tc c)}.
\end{equation}
\Cref{eq:nat-A-rhs} implies \cref{eq:nat-A-lhs} with $c_1 = c$.
Conversely, if \cref{eq:nat-A-lhs} holds, then $a \tc c \leq a \tc c_1$, so $\rel{(a \tc c)}{M}{(b \tc c)}$ by \cref{eq:closed-source}, which gives \cref{eq:nat-A-rhs}.

\paragraph{Naturality in $B$.}
Let $M \colon A \otimes C \profto B \otimes C$ and $N \colon B \profto B'$.
Then $\tup{a, b'}$ is an element of $\Tr^{C}_{A,B'}(M \then (N \otimes \id_C))$ if and only if
\begin{equation}
    \label{eq:nat-B-lhs}
    \text{there are } c, c_1 \in C \text{ and } b \in B \text{ with } \rel{(a \tc c)}{M}{(b \tc c_1)}, \quad \rel{b}{N}{b'}, \quad c_1 \leq c,
\end{equation}
and it is an element of $\Tr^{C}_{A,B}(M) \then N$ if and only if
\begin{equation}
    \label{eq:nat-B-rhs}
    \text{there are } c \in C \text{ and } b \in B \text{ with } \rel{(a \tc c)}{M}{(b \tc c)}, \quad \rel{b}{N}{b'}.
\end{equation}
\Cref{eq:nat-B-rhs} implies \cref{eq:nat-B-lhs} with $c_1 = c$.
Conversely, if \cref{eq:nat-B-lhs} holds, then $b \tc c_1 \leq b \tc c$, so $\rel{(a \tc c)}{M}{(b \tc c)}$ by \cref{eq:closed-target}, which gives \cref{eq:nat-B-rhs}.

\paragraph{Dinaturality in $C$, for all morphisms.}
Let $M \colon A \otimes C \profto B \otimes C'$ and $N \colon C' \profto C$.
The pair $\tup{a,b}$ is an element of $\Tr^{C}_{A,B}(M \then (\id_B \otimes N))$ if and only if there are $c \in C$, $c' \in C'$, and $b_1 \in B$ with
\begin{equation*}
    \rel{(a \tc c)}{M}{(b_1 \tc c')}, \qquad b_1 \leq b, \qquad \rel{c'}{N}{c}.
\end{equation*}
By \cref{eq:closed-target}, applied to $b_1 \tc c' \leq b \tc c'$, this holds if and only if
\begin{equation}
    \label{eq:dinat-middle}
    \text{there are } c \in C \text{ and } c' \in C' \text{ with } \rel{(a \tc c)}{M}{(b \tc c')} \text{ and } \rel{c'}{N}{c}.
\end{equation}
The pair $\tup{a,b}$ is an element of $\Tr^{C'}_{A,B}((\id_A \otimes N) \then M)$ if and only if there are $c' \in C'$, $c \in C$, and $a_1 \in A$ with
\begin{equation*}
    a \leq a_1, \qquad \rel{c'}{N}{c}, \qquad \rel{(a_1 \tc c)}{M}{(b \tc c')}.
\end{equation*}
By \cref{eq:closed-source}, applied to $a \tc c \leq a_1 \tc c$, this also holds if and only if \cref{eq:dinat-middle} holds.

\paragraph{Vanishing I.}
Let $M \colon A \profto B$; recall that $A \otimes \tL{\,} = A$ and $B \otimes \tL{\,} = B$.
The only element of $\tL{\,}$ is $\etup$, and $a \tc \etup = a$, so $\tup{a,b} \in \Tr^{\tL{\,}}_{A,B}(M)$ if and only if $\rel{a}{M}{b}$.

\paragraph{Vanishing II.}
Let $M \colon A \otimes C \otimes D \profto B \otimes C \otimes D$.
By definition, $\Tr^{D}_{A \otimes C, B \otimes C}(M)$ relates $a \tc c$ to $b \tc c'$ precisely when there is $d \in D$ with $\rel{(a \tc c \tc d)}{M}{(b \tc c' \tc d)}$.
So $\tup{a,b}$ is an element of $\Tr^{C}_{A,B}(\Tr^{D}_{A \otimes C, B \otimes C}(M))$ if and only if there are $c \in C$ and $d \in D$ with $\rel{(a \tc c \tc d)}{M}{(b \tc c \tc d)}$.
Since the elements of $C \otimes D$ are precisely the tuples $c \tc d$, this is the condition for $\tup{a,b}$ to be an element of $\Tr^{C \otimes D}_{A,B}(M)$.

\paragraph{Superposing.}
Let $M \colon A \otimes C \profto B \otimes C$ and $N \colon E \profto F$.
Then $N \otimes M$ relates $e \tc a \tc c$ to $f \tc b \tc c'$ precisely when $\rel{e}{N}{f}$ and $\rel{(a \tc c)}{M}{(b \tc c')}$.
So $\Tr^{C}_{E \otimes A, F \otimes B}(N \otimes M)$ relates $e \tc a$ to $f \tc b$ precisely when $\rel{e}{N}{f}$ and there is $c \in C$ with $\rel{(a \tc c)}{M}{(b \tc c)}$, that is, precisely when $N \otimes \Tr^{C}_{A,B}(M)$ does.

\paragraph{Yanking.}
By \cref{eq:mRelL-braiding}, $\sigma_{C,C}$ relates $c \tc c''$ to $c' \tc c''$ precisely when $c \leq c''$ and $c'' \leq c'$.
So $\tup{c,c'} \in \Tr^{C}_{C,C}(\sigma_{C,C})$ if and only if there is $c''$ with $c \leq c'' \leq c'$, that is, if and only if $c \leq c'$.
Hence $\Tr^{C}_{C,C}(\sigma_{C,C}) = {\leq_C} = \id_C$.
This completes the proof of \cref{prop:mRelL-trace}. \qed

\begin{thebibliography}{9}
    \bibitem{dickson-survey}
    E. Aichinger and F. Aichinger,
    \emph{Dickson's lemma, Higman's theorem and beyond: a survey of some basic results in order theory},
    Expositiones Mathematicae (2020), arXiv:1812.03024.

    \bibitem{bkpv}
    M. B\'{\i}lkov\'a, A. Kurz, D. Petri\c{s}an, and J. Velebil,
    \emph{Relation liftings on preorders and posets},
    arXiv:1210.1433.

    \bibitem{censi}
    A. Censi,
    \emph{A mathematical theory of co-design},
    arXiv:1512.08055.

    \bibitem{fong-spivak}
    B. Fong and D. I. Spivak,
    \emph{An Invitation to Applied Category Theory: Seven Sketches in Compositionality},
    Cambridge University Press, 2019.

    \bibitem{higman}
    G. Higman,
    \emph{Ordering by divisibility in abstract algebras},
    Proceedings of the London Mathematical Society (3) 2 (1952), 326--336.

    \bibitem{kock}
    A. Kock,
    \emph{Strong functors and monoidal monads},
    Archiv der Mathematik 23 (1972), 113--120.

    \bibitem{power-robinson}
    J. Power and E. Robinson,
    \emph{Premonoidal categories and notions of computation},
    Mathematical Structures in Computer Science 7 (1997), 453--468.
\end{thebibliography}

\end{document}
